\documentclass[a4paper,14pt]{extarticle}

\usepackage{fontspec}
\usepackage{polyglossia}
\setmainlanguage{russian}
\setotherlanguage{english}
\setmainfont[Ligatures=TeX]{Times New Roman}
\newfontfamily\cyrillicfont[Ligatures=TeX]{Times New Roman}
\newfontfamily\cyrillicfonttt{Menlo}[Scale=0.88]
\newfontfamily\cyrillicfontsf{Arial}

\usepackage[
  left=30mm,
  right=15mm,
  top=20mm,
  bottom=20mm
]{geometry}
\usepackage{setspace}
\onehalfspacing
\setlength{\parindent}{1.25cm}
\setlength{\parskip}{0pt}

% Сдаваемая версия ВКР не должна разрывать слова переносами вида
% "обозначе- / ния", особенно на стыке страниц.
\hyphenpenalty=10000
\exhyphenpenalty=10000
\brokenpenalty=10000
\doublehyphendemerits=1000000
\finalhyphendemerits=1000000

\usepackage{indentfirst}
\usepackage{microtype}
\usepackage{graphicx}
\usepackage{adjustbox}
\usepackage{float}
\usepackage{flafter}
\usepackage{caption}
\usepackage{longtable}
\usepackage{booktabs}
\usepackage{array}
\usepackage{calc}
\usepackage{ragged2e}
\usepackage{etoolbox}
\usepackage{enumitem}
\usepackage{amsmath}
\usepackage{amssymb}
\usepackage{fvextra}
\usepackage{listings}
\usepackage{lastpage}
\usepackage{titlesec}
\usepackage{tocloft}
\usepackage{needspace}
\usepackage{placeins}
\usepackage{xurl}
\usepackage{hyperref}
\usepackage{bookmark}
\usepackage{xcolor}

\graphicspath{{./}{../}{docx_pipeline/figures/png/}{docx_pipeline/figures/svg/}}

\hypersetup{
  unicode=true,
  colorlinks=false,
  pdfborder={0 0 0},
  pdftitle={Разработка и внедрение интеллектуальной системы управления микроклиматом в теплице на основе машинного обучения},
  pdfauthor={Ларионов Сергей Иванович}
}

\setcounter{secnumdepth}{0}
\setcounter{tocdepth}{3}

\titleformat{\section}
  {\normalfont\bfseries}
  {}
  {0pt}
  {}
\titlespacing*{\section}{1.25cm}{0pt}{12pt}

\titleformat{\subsection}
  {\normalfont\bfseries}
  {}
  {0pt}
  {}
\titlespacing*{\subsection}{1.25cm}{12pt}{6pt}

\titleformat{\subsubsection}
  {\normalfont\bfseries}
  {}
  {0pt}
  {}
\titlespacing*{\subsubsection}{1.25cm}{10pt}{6pt}

\pretocmd{\section}{\FloatBarrier\clearpage}{}{}
\pretocmd{\subsection}{\Needspace{6\baselineskip}}{}{}
\pretocmd{\subsubsection}{\Needspace{5\baselineskip}}{}{}

\newcommand{\VkrStructuralSection}[1]{%
  \FloatBarrier
  \clearpage
  \phantomsection
  \addcontentsline{toc}{section}{#1}%
  \begin{center}\bfseries #1\end{center}%
  \par\nopagebreak\vspace{12pt}%
}

\newcommand{\VkrStructuralSectionNoToc}[1]{%
  \begin{center}\bfseries #1\end{center}%
  \par\nopagebreak\vspace{12pt}%
}

\newcommand{\VkrAppendixSection}[2]{%
  \FloatBarrier
  \clearpage
  \phantomsection
  \addcontentsline{toc}{section}{ПРИЛОЖЕНИЕ #1 #2}%
  \begin{center}\bfseries ПРИЛОЖЕНИЕ #1\end{center}%
  \par\nopagebreak\vspace{6pt}%
  \begin{center}\bfseries #2\end{center}%
  \par\nopagebreak\vspace{12pt}%
}

\renewcommand{\cftsecleader}{\cftdotfill{\cftdotsep}}
\makeatletter
\renewcommand{\tableofcontents}{%
  \section*{СОДЕРЖАНИЕ}%
  \@starttoc{toc}%
}
\makeatother
\setlength{\cftbeforesecskip}{2pt}
\setlength{\cftbeforesubsecskip}{1pt}

\captionsetup{
  font={normalsize,stretch=1},
  labelformat=empty,
  labelsep=none,
  justification=centering,
  singlelinecheck=false,
  skip=6pt
}
\captionsetup[lstlisting]{
  font={normalsize,stretch=1},
  labelformat=simple,
  labelsep=endash,
  justification=raggedright,
  singlelinecheck=false,
  skip=6pt
}

\setlist[itemize]{label=\textemdash,leftmargin=1.25cm,itemsep=0pt,topsep=0pt,parsep=0pt}
\setlist[enumerate]{leftmargin=1.25cm,itemsep=0pt,topsep=0pt,parsep=0pt}

\lstset{
  basicstyle=\ttfamily\small,
  breaklines=true,
  columns=fullflexible,
  keepspaces=true,
  frame=single,
  framerule=0.4pt,
  framesep=4pt,
  captionpos=t,
  rulecolor=\color{black},
  xleftmargin=0pt,
  xrightmargin=0pt,
  aboveskip=6pt,
  belowskip=6pt
}
\renewcommand{\lstlistingname}{Листинг}
\RecustomVerbatimEnvironment{verbatim}{Verbatim}{
  breaklines=true,
  breakanywhere=true,
  fontsize=\small,
  frame=single,
  framerule=0.4pt,
  framesep=4pt
}
\renewcommand{\texttt}[1]{\textnormal{#1}}

\sloppy
\emergencystretch=2em
\clubpenalty=10000
\widowpenalty=10000
\displaywidowpenalty=10000
\pagestyle{plain}

\providecommand{\tightlist}{%
  \setlength{\itemsep}{0pt}\setlength{\parskip}{0pt}}
\newcommand{\VkrFigureFull}[1]{%
  \adjustbox{max width=\linewidth,max totalheight=0.74\textheight,keepaspectratio,center}{#1}%
}
\newcommand{\VkrFigureInline}[1]{%
  \adjustbox{max width=0.94\linewidth,max totalheight=0.55\textheight,keepaspectratio,center}{#1}%
}
\newcommand{\VkrFigureArchitecture}[1]{%
  \adjustbox{max width=0.98\linewidth,max totalheight=0.74\textheight,keepaspectratio,center}{#1}%
}
\newcommand{\VkrFigureTurnCycle}[1]{%
  \adjustbox{max width=0.98\linewidth,max totalheight=0.76\textheight,keepaspectratio,center}{#1}%
}
\newcommand{\VkrFigureStagePolicy}[1]{%
  \adjustbox{max width=0.98\linewidth,max totalheight=0.65\textheight,keepaspectratio,center}{#1}%
}
\newcommand{\VkrFigureExperimentDesign}[1]{%
  \adjustbox{max width=0.98\linewidth,max totalheight=0.60\textheight,keepaspectratio,center}{#1}%
}
\newcommand{\VkrFigureTall}[1]{%
  \adjustbox{max width=0.88\linewidth,max totalheight=0.50\textheight,keepaspectratio,center}{#1}%
}
\newcommand{\VkrFigureVertical}[1]{%
  \adjustbox{max width=0.88\linewidth,max totalheight=0.56\textheight,keepaspectratio,center}{#1}%
}
\newcommand{\VkrFigureCompact}[1]{%
  \adjustbox{max width=0.92\linewidth,max totalheight=0.48\textheight,keepaspectratio,center}{#1}%
}
\newcommand{\VkrFigureWideCompact}[1]{%
  \adjustbox{max width=0.88\linewidth,max totalheight=0.38\textheight,keepaspectratio,center}{#1}%
}
\providecommand{\pandocbounded}[1]{\VkrFigureFull{#1}}
\providecommand{\passthrough}[1]{#1}
\newcounter{none}

\setlength{\LTpre}{0pt}
\setlength{\LTpost}{6pt}
\setlength{\tabcolsep}{3.5pt}
\setlength{\extrarowheight}{2pt}
\renewcommand{\arraystretch}{1.22}
\AtBeginEnvironment{longtable}{\small\singlespacing}

\newcommand{\VkrTableCaption}[1]{%
  \Needspace{0.44\textheight}%
  \par\smallskip
  {\singlespacing\noindent #1\par}%
  \vspace{2pt}%
}


\begin{document}
\hypersetup{pageanchor=false}
\newcommand{\VkrTitleHeaderLine}[1]{%
  \noindent{\fontsize{12pt}{14pt}\selectfont\makebox[\textwidth][c]{#1}\par}%
}
\newcommand{\VkrUdkField}{УДК \underline{\hspace{1.8cm}}}

\begin{titlepage}
\fontsize{12pt}{14pt}\selectfont
\singlespacing
\VkrTitleHeaderLine{МИНИСТЕРСТВО СЕЛЬСКОГО ХОЗЯЙСТВА РОССИЙСКОЙ ФЕДЕРАЦИИ}

\vspace{0.55cm}

\VkrTitleHeaderLine{ФЕДЕРАЛЬНОЕ ГОСУДАРСТВЕННОЕ БЮДЖЕТНОЕ}
\VkrTitleHeaderLine{ОБРАЗОВАТЕЛЬНОЕ УЧРЕЖДЕНИЕ ВЫСШЕГО ОБРАЗОВАНИЯ}
\VkrTitleHeaderLine{«САНКТ-ПЕТЕРБУРГСКИЙ ГОСУДАРСТВЕННЫЙ}
\VkrTitleHeaderLine{АГРАРНЫЙ УНИВЕРСИТЕТ»}

\vspace{0.75cm}

\VkrTitleHeaderLine{ИНСТИТУТ ЭНЕРГЕТИКИ}

\vspace{0.75cm}

\VkrTitleHeaderLine{КАФЕДРА ТЕПЛОТЕХНИКИ И ЭНЕРГООБЕСПЕЧЕНИЯ}

\vspace{0.9cm}

\begin{flushright}
На правах рукописи

\vspace{0.2cm}

\VkrUdkField
\end{flushright}

\vspace{0.75cm}

\begin{center}
Ларионов Сергей Иванович

\vspace{0.8cm}

\textbf{«Разработка и внедрение интеллектуальной системы управления микроклиматом в теплице на основе машинного обучения для повышения эффективности выращивания растений»}

\vspace{0.85cm}

Выпускная квалификационная работа магистра

\vspace{0.65cm}

Направление подготовки\\
«Теплотехника и теплоэнергетика в агропромышленном комплексе»
\end{center}

\vspace{0.4cm}

\begin{flushright}
\begin{minipage}{0.46\textwidth}
Выпускная квалификационная\\
работа защищена

\vspace{0.15cm}

«\underline{\hspace{0.8cm}}»\underline{\hspace{2.5cm}}2026 г.

\vspace{0.18cm}

Оценка

\vspace{0.12cm}

\underline{\hspace{3.0cm}}

\vspace{0.18cm}

Секретарь ГЭК \underline{\hspace{2.2cm}}
\end{minipage}
\end{flushright}

\vfill

\begin{center}
г. Санкт-Петербург\\
2026
\end{center}
\end{titlepage}

\setcounter{page}{2}
\hypersetup{pageanchor=true}

\VkrStructuralSectionNoToc{РЕФЕРАТ}\label{ux440ux435ux444ux435ux440ux430ux442}

Работа содержит \pageref{LastPage} с., 0 рис., 6 табл., 32 источн.

Выпускная квалификационная работа посвящена разработке интеллектуальной
системы управления микроклиматом теплицы на основе машинного обучения.
Теплица рассматривается как теплотехнический и киберфизический объект, в
котором температурный режим, влажность, воздухообмен и концентрация
углекислого газа должны анализироваться совместно с управляющими
воздействиями.

Цель работы состоит в обосновании и разработке расширяемой
IoT/ML-платформы для мониторинга, журналирования и последующего
прогнозного управления микроклиматом теплицы. В текущей редакции
основное внимание уделено постановке задачи, анализу литературы и
обоснованию архитектурной рамки, в которой разнородные датчики и
исполнительные устройства подключаются через единую модель обмена
данными.

В аналитической части рассмотрены современные подходы к smart
greenhouse, IoT-архитектурам, прогнозированию микроклимата,
предиктивному управлению, обучению с подкреплением, энергоэффективности
и обработке неоднородных потоков сенсорных данных. На основе обзора
сформулированы требования к системе как к расширяемой платформе,
способной не только фиксировать параметры микроклимата, но и накапливать
данные для дальнейшего применения методов машинного обучения.

Ключевые слова: теплица, микроклимат, тепловлажностный режим,
интеллектуальное управление, интернет вещей, машинное обучение,
прогнозирование, энергоэффективность.

\clearpage

\VkrStructuralSectionNoToc{СОДЕРЖАНИЕ}
\makeatletter\@starttoc{toc}\makeatother
\clearpage

\VkrStructuralSection{ПЕРЕЧЕНЬ СОКРАЩЕНИЙ И ОБОЗНАЧЕНИЙ}\label{ux441ux43fux438ux441ux43eux43a-ux441ux43eux43aux440ux430ux449ux435ux43dux438ux439-ux438-ux43eux431ux43eux437ux43dux430ux447ux435ux43dux438ux439}

{\def\LTcaptype{none} % do not increment counter
\begin{longtable}[]{@{}
  >{\RaggedRight\arraybackslash}p{(\linewidth - 4\tabcolsep) * \real{0.3333}}
  >{\RaggedRight\arraybackslash}p{(\linewidth - 4\tabcolsep) * \real{0.3333}}
  >{\RaggedRight\arraybackslash}p{(\linewidth - 4\tabcolsep) * \real{0.3333}}@{}}
\hline
\begin{minipage}[b]{\linewidth}\RaggedRight
Сокращение
\end{minipage} & \begin{minipage}[b]{\linewidth}\RaggedRight
Расшифровка
\end{minipage} & \begin{minipage}[b]{\linewidth}\RaggedRight
Пояснение в контексте работы
\end{minipage} \\
\hline
\endhead
\hline
\endlastfoot
АСУ & автоматизированная система управления & система наблюдения и
управления параметрами микроклимата \\
ВКР & выпускная квалификационная работа & итоговая квалификационная
работа магистранта \\
ИИ & искусственный интеллект & методы анализа данных и поддержки
принятия решений \\
IoT & internet of things & подход к объединению датчиков, исполнительных
устройств и программных сервисов \\
ML & machine learning & машинное обучение, применяемое для анализа и
прогнозирования временных рядов \\
MQTT & message queuing telemetry transport & легкий протокол обмена
сообщениями в IoT-системах \\
CO\(_2\) & carbon dioxide & концентрация углекислого газа в воздухе
теплицы \\
MPC & model predictive control & модельно-предиктивное управление \\
RL & reinforcement learning & обучение с подкреплением \\
MAE & mean absolute error & средняя абсолютная ошибка прогноза \\
RMSE & root mean squared error & корень из средней квадратичной ошибки
прогноза \\
\end{longtable}
}

\VkrStructuralSection{ВВЕДЕНИЕ}\label{ux432ux432ux435ux434ux435ux43dux438ux435}

Современные тепличные комплексы требуют точного контроля параметров
микроклимата: температуры, влажности, концентрации CO\(_2\), режимов
вентиляции, нагрева, увлажнения и полива. С теплотехнической точки
зрения теплица является объектом с тепловым балансом, влажностным
режимом и вентиляционным воздухообменом. Даже в малой теплице параметры
изменяются не изолированно: включение вентилятора влияет на температуру
и влажность, открытие заслонок меняет воздухообмен, работа нагревателей
имеет инерционный эффект, а концентрация CO\(_2\) зависит от вентиляции,
фотосинтетической активности растений и притока наружного воздуха.
Поэтому управление микроклиматом нельзя сводить только к ручному
включению оборудования или простым пороговым действиям без накопления
истории.

Актуальность работы связана с переходом тепличных систем от локальной
автоматики к киберфизическим IoT-решениям. В современных исследованиях
smart greenhouse рассматривается как система, объединяющая датчики,
исполнительные механизмы, сетевой обмен, хранилище данных, веб-интерфейс
и интеллектуальный аналитический слой {[}1; 2{]}. Такая постановка
особенно важна для малых и учебно-экспериментальных теплиц: если система
проектируется только под один датчик и одно реле, она быстро перестает
быть исследовательской платформой. Для дальнейшего применения машинного
обучения требуется расширяемая архитектура, в которой новые сенсорные
узлы, исполнительные каналы и расчетные признаки подключаются без
переписывания всей системы.

Отдельное значение имеет энергоэффективность. Нагрев, вентиляция,
увлажнение и освещение потребляют значительные ресурсы, а вентиляционные
воздействия могут приводить к теплопотерям {[}28; 31{]}. При этом
реакция теплицы на управляющее воздействие зависит от текущего состояния
среды, внешних условий, инерции конструкций, плотности растений и режима
воздухообмена. Поэтому управление должно учитывать не только текущие
значения датчиков, но и прогноз развития микроклимата {[}14; 25{]}.

Методы машинного обучения позволяют использовать накопленную телеметрию
для прогнозирования температуры, влажности, CO\(_2\), роста растений и
урожайности. В литературе применяются LSTM, RNN, TCN, Transformer и
вероятностные модели, ориентированные на временные ряды тепличных данных
{[}9; 10; 11; 13{]}. Однако для практического применения таких моделей
сначала необходим надежный инженерный слой: датчики должны стабильно
публиковать данные, исполнительные устройства должны управляться через
понятный интерфейс, а история измерений и воздействий должна сохраняться
в пригодном для анализа виде.

В рамках данной работы рассматривается разработка и опытное внедрение
расширяемой программно-аппаратной IoT/ML-платформы мониторинга и
управления микроклиматом теплицы. Платформа объединяет микроконтроллеры,
MQTT-брокер, веб-панель оператора, правила автоматизации, профили
управления, журналирование данных и подготовку истории для машинного
обучения. Реальная история измерений и управляющих воздействий
необходима для перехода от эмпирического управления к имитационной или
ML-модели тепловлажностного режима {[}32; 29{]}. Машинное обучение в
работе рассматривается как верхний интеллектуальный слой, который
опирается на данные, собранные инженерной системой. Такой подход
позволяет сохранить надежность базового управления и одновременно
подготовить основу для прогнозных моделей и дальнейшего
интеллектуального управления.

Цель работы: разработать и внедрить расширяемую IoT/ML-платформу
мониторинга и управления тепловлажностным режимом теплицы с возможностью
подключения разнородных датчиков, управления исполнительными
устройствами, накопления истории и последующего применения методов
машинного обучения для прогнозирования и интеллектуальной поддержки
управления.

Для достижения цели необходимо решить следующие задачи:

\begin{enumerate}
\def\labelenumi{\arabic{enumi})}
\tightlist
\item
  Проанализировать современные подходы к автоматизации микроклимата
  теплиц, IoT-архитектурам, прогнозированию микроклимата и
  интеллектуальному управлению.
\item
  Определить требования к расширяемой системе мониторинга и управления:
  состав параметров, датчики, исполнительные механизмы, канал связи,
  хранение данных, интерфейс оператора и пригодность истории для
  ML-анализа.
\item
  Спроектировать архитектуру системы на базе микроконтроллеров, MQTT,
  серверного приложения, базы данных истории и единой модели обмена
  сообщениями.
\item
  Реализовать программно-аппаратный контур: сенсорные узлы, релейный
  контроллер, MQTT-обмен, dashboard, журналирование показаний и
  состояний.
\item
  Подготовить основу для интеллектуального слоя: накопление временных
  рядов, экспорт данных, правила и профили управления, план
  ML-эксперимента.
\item
  Проверить работоспособность системы на опытном объекте, сформировать
  датасет эксплуатации и определить направления дальнейшего развития.
\end{enumerate}

Объект исследования: тепловлажностный режим теплицы как управляемая
киберфизическая система агропромышленного комплекса.

Предмет исследования: методы и программно-аппаратные средства
мониторинга, управления и интеллектуальной поддержки регулирования
микроклимата теплицы.

Практическая значимость работы состоит в создании системы, которая может
использоваться как основа для эксплуатации малой теплицы и для
дальнейших экспериментов с прогнозированием микроклимата. Реализованная
архитектура позволяет собирать телеметрию, управлять исполнительными
устройствами, хранить историю и формировать датасет для машинного
обучения.

Научная новизна работы заключается в архитектурно-методическом подходе к
построению малой теплицы как расширяемой IoT/ML-платформы. В отличие от
одноразовой схемы «датчик - реле - панель управления», в работе
рассматривается единый контур, где разнородные датчики публикуют данные
в общей модели обмена, управляющие воздействия регистрируются совместно
с параметрами микроклимата, а накопленная история становится основой для
последующего ML-прогнозирования реакции теплицы. При этом не заявляется
разработка нового алгоритма машинного обучения: интеллектуальный слой
рассматривается как прогнозная надстройка над безопасным правиловым
контуром.

Структура работы включает введение, три главы, заключение, список
литературы и приложения. В первой главе приведен аналитический обзор
источников по smart greenhouse, IoT, прогнозированию и методам
управления. Во второй главе описано проектирование системы. В третьей
главе рассмотрены реализация, проверка работоспособности и подготовка
ML-эксперимента.

\section{1 Аналитический обзор систем интеллектуального управления
микроклиматом
теплицы}\label{ux430ux43dux430ux43bux438ux442ux438ux447ux435ux441ux43aux438ux439-ux43eux431ux437ux43eux440-ux441ux438ux441ux442ux435ux43c-ux438ux43dux442ux435ux43bux43bux435ux43aux442ux443ux430ux43bux44cux43dux43eux433ux43e-ux443ux43fux440ux430ux432ux43bux435ux43dux438ux44f-ux43cux438ux43aux440ux43eux43aux43bux438ux43cux430ux442ux43eux43c-ux442ux435ux43fux43bux438ux446ux44b}

\subsection{1.1 Теплица как киберфизическая система
управления}\label{ux442ux435ux43fux43bux438ux446ux430-ux43aux430ux43a-ux43aux438ux431ux435ux440ux444ux438ux437ux438ux447ux435ux441ux43aux430ux44f-ux441ux438ux441ux442ux435ux43cux430-ux443ux43fux440ux430ux432ux43bux435ux43dux438ux44f}

Современная теплица является не просто защищенным сооружением для
выращивания растений, а киберфизической системой, в которой датчики,
исполнительные механизмы, микроконтроллеры, серверное программное
обеспечение и алгоритмы принятия решений работают как единый контур.
Целевыми параметрами такого контура выступают температура воздуха,
относительная влажность, концентрация CO\(_2\), освещенность, влажность
почвы или субстрата, режим полива и состояние растений. Поддержание этих
параметров влияет на урожайность, скорость роста, устойчивость растений
к стрессу и энергоемкость выращивания.

В обзорной литературе последних лет тепличное хозяйство все чаще
рассматривается через призму smart greenhouse, то есть управляемой
среды, где IoT, искусственный интеллект, спектральные датчики,
мультимодальное слияние данных и интеллектуальные системы используются
для повышения эффективности выращивания. В обзоре El Ouaham et
al.~показано, что развитие smart greenhouse связано не только с заменой
ручного контроля на датчики, но и с переходом к системам, которые
собирают разнородные данные, анализируют их и поддерживают принятие
решений {[}1{]}. Для данной ВКР это принципиально: ценность системы
определяется не только тем, что она измеряет температуру и включает
реле, а тем, что она создает расширяемый контур данных для последующего
прогнозирования и анализа реакции теплицы.

Русскоязычные источники подтверждают ту же постановку на инженерном
уровне. В сравнительном анализе современных систем автоматизированного
управления микроклиматом подчеркивается, что автоматизация теплиц
направлена на снижение эксплуатационных затрат, повышение урожайности и
переход от простого поддержания параметров к интеллектуальному
управлению {[}3{]}. Халина и Цыбанюк описывают базовую структуру
автоматизированной системы: датчики измеряют параметры среды, блок
управления сравнивает их с заданными значениями, после чего
исполнительные механизмы изменяют температуру, влажность, освещение или
другие параметры {[}4{]}. Такая схема является исходной, но в
современных условиях она требует расширения прогнозными и аналитическими
функциями.

Особенность тепличного микроклимата состоит в инерционности и
взаимосвязанности параметров. Нагрев, вентиляция, увлажнение, подача
CO\(_2\) и освещение влияют друг на друга, а реакция растений
проявляется не мгновенно. Тарков и Сукьясов рассматривают задачу
управления параметрами микроклимата через уравнения баланса и отмечают
проблему запаздывания процесса регулирования {[}5{]}. Это означает, что
реактивное управление по текущему отклонению от уставки не всегда
обеспечивает оптимальный результат: система может опоздать с
воздействием или создать избыточные энергетические затраты.

\subsection{1.2 Тепловлажностный режим и энергетическая рамка
теплицы}\label{ux442ux435ux43fux43bux43eux432ux43bux430ux436ux43dux43eux441ux442ux43dux44bux439-ux440ux435ux436ux438ux43c-ux438-ux44dux43dux435ux440ux433ux435ux442ux438ux447ux435ux441ux43aux430ux44f-ux440ux430ux43cux43aux430-ux442ux435ux43fux43bux438ux446ux44b}

Для направления «Теплотехника и теплоэнергетика в агропромышленном
комплексе» принципиально важно рассматривать теплицу как объект тепло- и
массообмена. Температура, влажность и CO\(_2\) формируются под действием
солнечной радиации, теплопотерь через ограждения, вентиляционного
воздухообмена, испарения и транспирации растений, нагрева, увлажнения и
работы вспомогательного оборудования. В обзоре Chen et al.~тепличная
система прямо описывается как сложная многовходовая и многовыходовая
система, где управление температурой, влажностью, светом, CO\(_2\) и
воздушными потоками должно одновременно обеспечивать рост растений и
энергосбережение {[}28{]}.

Динамические модели микроклимата теплицы обычно опираются на баланс
энергии, энтальпии и массы. Это важно для текущей ВКР: даже если
практическая реализация использует датчики и реле, физическая сущность
задачи остается теплотехнической. Самарин и др. формулируют задачу
имитационного моделирования оптимального управления микроклиматом
сельскохозяйственного объекта через температурно-влажностный режим,
тепловой баланс, влажностное состояние конструкций и энергосберегающий
режим выращивания {[}32{]}. Такой подход поддерживает выбранную рамку:
система должна не только включать устройства, но и формировать данные
для модели поведения теплицы.

Отдельная проблема состоит в пространственной неоднородности
микроклимата. Ogunlowo et al.~показывают, что распределение тепла и
массы в одно- и многопролетных теплицах влияет на точность оценки
энергетической потребности; при игнорировании распределения параметров
возможна переоценка или недооценка расчетов {[}29{]}. Это обосновывает
применение нескольких измерительных узлов внутри теплицы и последующий
анализ расхождений между датчиками, а не только использование одного
среднего значения.

Вентиляция является одним из ключевых каналов воздействия на
микроклимат. Она одновременно снижает температуру, удаляет влагу,
изменяет воздухообмен и влияет на концентрацию CO\(_2\). Li et al.~на
данных девяти односкатных теплиц показывают, что открытие вентиляции
меняет температуру, относительную влажность, VPD и CO\(_2\), а также
зависит от плотности посадки, высоты растений, внутренних покрытий и
времени открытия вентиляционного проема {[}30{]}. Поэтому события
включения вентиляторов и приточных каналов в собственной системе должны
журналироваться вместе с температурой, влажностью и CO\(_2\): без этого
невозможно корректно анализировать реакцию теплицы на управляющее
воздействие.

Энергетический аспект особенно заметен при управлении вентиляцией и
отоплением. Ghaderi et al.~отмечают, что вентиляционные потери тепла
являются одним из значимых факторов энергетической эффективности теплиц,
а климат-контроль может составлять основную долю энергопотребления
тепличного хозяйства {[}31{]}. Следовательно, интеллектуальный слой
управления должен стремиться не только удерживать параметры в допустимом
диапазоне, но и снижать лишние переключения, преждевременные теплопотери
и компенсирующие действия после запаздывания.

\subsection{1.3 IoT-архитектура и расширяемость тепличных
систем}\label{iot-ux430ux440ux445ux438ux442ux435ux43aux442ux443ux440ux430-ux438-ux440ux430ux441ux448ux438ux440ux44fux435ux43cux43eux441ux442ux44c-ux442ux435ux43fux43bux438ux447ux43dux44bux445-ux441ux438ux441ux442ux435ux43c}

Практическое внедрение интеллектуального управления начинается с
надежного сбора телеметрии. IoT-подходы позволяют связать датчики,
исполнительные устройства, контроллеры и программные сервисы в единую
инфраструктуру. Bersani et al.~рассматривают smart greenhouse в
контексте Industry 4.0 и показывают, что IoT-системы для теплиц включают
сенсорные сети, каналы связи, удаленный мониторинг и управляющие контуры
{[}2{]}. Такой подход хорошо согласуется с текущим проектом, где уже
присутствуют прошивки микроконтроллеров и серверный dashboard.

В инженерных реализациях часто используются микроконтроллеры, Raspberry
Pi, Arduino/ESP, MQTT, веб-интерфейсы и базы данных. Кельсин в ВКР по
автоматизированной системе управления тепличным хозяйством рассматривает
архитектуру на микроконтроллерах, близкую к практическим учебным и малым
производственным системам {[}7{]}. Статья по MQTT-based smart greenhouse
monitoring дополнительно показывает, почему MQTT удобен для телеметрии:
это легкий протокол обмена сообщениями, пригодный для IoT-устройств с
ограниченными ресурсами {[}23{]}. Для текущей ВКР это важно, поскольку
коммуникационный слой должен быть достаточно простым для
микроконтроллеров и достаточно надежным для накопления данных.

Для перехода от локальной автоматики к платформе недостаточно просто
добавить веб-интерфейс. Система должна иметь расширяемую модель
подключаемых устройств: разные датчики могут измерять температуру,
влажность, CO\(_2\), освещенность, влажность субстрата или
дополнительные расчетные параметры, но для верхнего уровня они должны
попадать в историю в сопоставимой структуре. Такой подход согласуется с
направлением multimodal data fusion в smart greenhouse: данные от разных
источников становятся не разрозненными показаниями, а общей
информационной базой для анализа состояния растений и микроклимата
{[}1{]}. В практической архитектуре это означает необходимость
идентификатора устройства, типа метрики, времени измерения, единиц
измерения и признака качества данных.

Расширяемость важна и для исполнительных устройств. В малой теплице реле
может управлять вентилятором, заслонкой, насосом, увлажнителем или
нагревателем, но для анализа микроклимата все эти действия должны
фиксироваться как события воздействия на тепловлажностный режим. Без
совместной регистрации датчиков и реле невозможно отличить естественное
изменение температуры от реакции на вентиляцию или нагрев. Поэтому
журнал исполнительных воздействий является не вспомогательной функцией
dashboard, а обязательной частью ML-ready датасета.

Прикладные smart farming эксперименты демонстрируют, что тепличная
система должна рассматриваться как полный стек: датчики, исполнительные
механизмы, логика управления, хранение данных и интерфейс оператора.
Joni et al.~описывают IoT-enhanced greenhouse design для выращивания
риса, где IoT-инфраструктура связывает параметры среды и агротехническую
задачу {[}24{]}. Несмотря на отличие культуры, работа полезна как пример
построения комплексной системы. Thwin et al.~идут еще дальше и связывают
прогноз микроклимата с web integration и adaptive equipment control
{[}16{]}. Это особенно важно для данной ВКР, так как собственный проект
уже содержит dashboard, прошивки и историю событий, а значит может быть
описан как первый инкремент интеллектуального контура, а не как
завершенная одноцелевая автоматика.

Отдельно следует отметить работы, где IoT-инфраструктура связывается с
reinforcement learning. Platero-Horcajadas et al.~рассматривают
интеграцию IoT и RL для оптимизированного climate control {[}21{]}. Для
текущей ВКР этот источник важен не как требование немедленно реализовать
RL-агента, а как подтверждение архитектурного тезиса: интеллектуальный
регулятор возможен только там, где есть непрерывная телеметрия,
наблюдаемые исполнительные воздействия и воспроизводимый обмен данными.

\subsection{1.4 Прогнозирование микроклимата, роста и
урожайности}\label{ux43fux440ux43eux433ux43dux43eux437ux438ux440ux43eux432ux430ux43dux438ux435-ux43cux438ux43aux440ux43eux43aux43bux438ux43cux430ux442ux430-ux440ux43eux441ux442ux430-ux438-ux443ux440ux43eux436ux430ux439ux43dux43eux441ux442ux438}

Накопление телеметрии само по себе не делает систему интеллектуальной.
Интеллектуальный слой появляется тогда, когда исторические данные
используются для прогнозирования будущего состояния микроклимата, роста
растений или урожайности. Alhnaity et al.~применяют LSTM для
прогнозирования роста и урожайности в greenhouse environments и
показывают, что последовательностные модели способны извлекать полезные
зависимости из временных рядов тепличных данных {[}9{]}. Gong et
al.~предлагают связку TCN и RNN для прогнозирования урожайности и
сравнивают ее с классическими ML-подходами {[}10{]}. Эти работы
обосновывают использование истории наблюдений, а не только текущего
значения датчика.

Для задачи управления микроклиматом особенно важны работы, где
прогнозируются непосредственно температура, влажность и CO\(_2\). Ahn et
al.~сравнивают transformer-based и RNN-based модели для предсказания
температуры, относительной влажности и концентрации CO\(_2\) в теплице
{[}11{]}. Huang et al.~используют Attention CNN-LSTM для прогнозирования
среды грибной теплицы, также работая с температурой, влажностью и
CO\(_2\) {[}12{]}. Эти источники показывают, что микроклимат является
естественным объектом для time-series prediction, а не только для
порогового контроля. При этом качество прогноза зависит не только от
выбора модели, но и от состава признаков: полезными могут быть лаги
датчиков, время суток, сезонность, состояние вентиляции, нагрева и
увлажнения, а также различия между точками измерения.

Отдельное направление связано с прогнозным управлением и
энергоэффективностью. Aborujilah et al.~рассматривают forecast-driven
climate control для smart greenhouse и используют LSTM-модель для
proactive adjustments, направленных на снижение энергетических затрат
{[}14{]}. Soumo et al.~описывают data-driven greenhouse climate
regulation при выращивании салата с использованием BiLSTM и GRU
predictive control {[}15{]}. Thwin et al.~предлагают multivariate
multistep LSTM, который прогнозирует микроклимат и взаимодействует с
системой управления оборудованием через веб-интеграцию {[}16{]}. В
совокупности эти работы позволяют обосновать, что прогнозный слой может
быть включен в практическую АСУ теплицы.

При этом прогнозы не следует рассматривать как абсолютно надежные. Choi
and Yang предлагают probabilistic deep learning framework для greenhouse
microclimate prediction с учетом time-varying uncertainty и covariance
analysis {[}13{]}. Для диплома этот источник важен методологически: даже
если в прототипе применяется более простая модель, в тексте необходимо
признавать неопределенность прогноза и необходимость защитных
ограничений нижнего уровня управления.

\subsection{1.5 Методы управления: PID, fuzzy, MPC, RL и эвристическая
оптимизация}\label{ux43cux435ux442ux43eux434ux44b-ux443ux43fux440ux430ux432ux43bux435ux43dux438ux44f-pid-fuzzy-mpc-rl-ux438-ux44dux432ux440ux438ux441ux442ux438ux447ux435ux441ux43aux430ux44f-ux43eux43fux442ux438ux43cux438ux437ux430ux446ux438ux44f}

Классические системы управления теплицей часто строятся на пороговых
правилах, PID-регуляторах или нечеткой логике. Их преимущество состоит в
простоте, интерпретируемости и возможности реализации на
микроконтроллерах. Воробьева и др. рассматривают интеллектуальную
систему контроля температуры в теплице и сравнивают PID-регулирование с
нечетким регулятором {[}6{]}. Для инженерного проекта такие методы
остаются важными, потому что нижний уровень управления должен быть
предсказуемым и безопасным.

Однако при усложнении задачи возникает потребность в методах, которые
учитывают будущую динамику, ограничения и стоимость ресурсов. Model
Predictive Control является одним из основных подходов в научной
литературе по greenhouse climate control, поскольку он использует модель
системы для оптимизации управляющих воздействий на горизонте прогноза.
Mallick et al.~показывают, что MPC явно учитывает физические
ограничения, но сильно зависит от точности модели; для компенсации
неопределенности они предлагают связку MPC и reinforcement learning
{[}17{]}. Msaad et al.~также рассматривают RL-guided MPC для автономного
управления теплицей {[}18{]}.

Сравнение MPC и RL важно для корректного позиционирования ВКР. Morcego
et al.~рассматривают MPC как model-based подход, а RL как learning-based
подход к greenhouse climate control {[}19{]}. MPC лучше интерпретируется
и естественно работает с ограничениями, тогда как RL способен улучшать
стратегию на основе данных, но требует осторожности из-за рисков
обучения и безопасности. Поэтому для практической ВКР разумно описывать
RL/MPC не как обязательную реализацию в полном объеме, а как верхний
интеллектуальный слой, который может выдавать рекомендации или
корректировать уставки, тогда как нижний контур остается защищенным.

Перспективным направлением является включение оператора или агронома в
цикл обучения. Xiao et al.~предлагают grower-in-the-loop interactive
reinforcement learning для greenhouse climate control {[}20{]}. Такой
подход хорошо соответствует практической логике внедрения: система может
не заменять оператора полностью, а поддерживать его решения и
использовать экспертную обратную связь.

Помимо MPC/RL, в литературе встречаются эвристические оптимизационные
методы. Jawad et al.~используют artificial bee colony algorithm для
energy optimization and plant comfort management в smart greenhouse
{[}22{]}. Этот источник показывает, что задача управления микроклиматом
может решаться не только нейросетями или MPC, но и другими
оптимизационными методами. Для текущей работы это полезно как обзор
альтернатив, но основная логика ВКР остается связанной с IoT,
прогнозированием и гибридным управлением.

\subsection{1.6 Энергоэффективность и устойчивость тепличного
управления}\label{ux44dux43dux435ux440ux433ux43eux44dux444ux444ux435ux43aux442ux438ux432ux43dux43eux441ux442ux44c-ux438-ux443ux441ux442ux43eux439ux447ux438ux432ux43eux441ux442ux44c-ux442ux435ux43fux43bux438ux447ux43dux43eux433ux43e-ux443ux43fux440ux430ux432ux43bux435ux43dux438ux44f}

Энергопотребление является одним из ключевых ограничений тепличного
производства. Отопление, вентиляция, увлажнение, освещение и подача
CO\(_2\) требуют ресурсов, а реактивное управление может приводить к
избыточным включениям исполнительных механизмов. Aborujilah et al.~прямо
указывают на проблему реактивных стратегий и предлагают forecast-driven
подход с LSTM для снижения энергетических потерь {[}14{]}. Эту же линию
поддерживает теплотехническая литература по вентиляционным потерям и
интеграции отопления, охлаждения и вентиляции {[}31{]}.

Более широкий энергетический контекст раскрывается в обзоре Wu and Fu по
Agricultural Energy Internet {[}25{]}. В нем greenhouse environmental
control рассматривается как энергетически управляемый процесс наряду с
ирригацией, дополнительным освещением, логистикой и
сельскохозяйственными машинами. Для ВКР этот источник полезен в разделе
актуальности: интеллектуальное управление микроклиматом не только
улучшает условия выращивания, но и связано с задачей рационального
использования энергии.

Энергоэффективность также связана с архитектурой управления. Если
система умеет прогнозировать состояние микроклимата, она может заранее
выбирать более мягкие воздействия, снижать частоту резких переключений и
избегать компенсирующих действий после запаздывания. Однако такие
преимущества возможны только при наличии надежной телеметрии,
устойчивого обмена данными и защитных правил нижнего уровня.

\subsection{1.7 Выводы по аналитическому
обзору}\label{ux432ux44bux432ux43eux434ux44b-ux43fux43e-ux430ux43dux430ux43bux438ux442ux438ux447ux435ux441ux43aux43eux43cux443-ux43eux431ux437ux43eux440ux443}

Проведенный обзор показывает, что тема интеллектуального управления
микроклиматом теплицы находится на пересечении нескольких направлений:
теплотехники, тепло- и массообмена, IoT-архитектур, автоматизированного
управления, машинного обучения, прогнозирования временных рядов,
энергоэффективности и человеко-ориентированного принятия решений.
Практические инженерные источники описывают датчики, контроллеры, MQTT,
dashboard и исполнительные механизмы, но часто ограничиваются пороговыми
или простыми регуляторами. Современные научные работы предлагают LSTM,
RNN, Transformer, Attention CNN-LSTM, вероятностные модели, MPC и RL,
однако не всегда доводят эти методы до конкретной малой теплицы с
реальной прошивкой, пользовательским интерфейсом и совместным
журналированием измерений и воздействий.

На этом фоне текущая ВКР может быть обоснована как разработка и опытное
внедрение расширяемой IoT/ML-платформы управления микроклиматом, которая
создает основу для интеллектуального слоя. Нижний уровень системы должен
обеспечивать сбор данных, безопасное управление и связь с
исполнительными устройствами. Верхний уровень должен накапливать
телеметрию, анализировать временные ряды, формировать прогнозы и
поддерживать корректировку уставок. Научный разрыв, закрываемый работой,
состоит не в создании нового ML-алгоритма, а в переходе от разрозненной
автоматики к воспроизводимой платформе: разнородные датчики подключаются
через единую модель обмена, события реле регистрируются вместе с
микроклиматом, а история становится пригодной для ML-прогнозирования
реакции теплицы. Такой гибридный подход согласует надежность
практической АСУ с современными направлениями ML/MPC/RL и позволяет
постепенно развивать систему без риска нереалистичного обещания
полностью автономного AI-управления.

\section{2 Проектирование расширяемой IoT/ML-платформы управления
микроклиматом
теплицы}\label{ux43fux440ux43eux435ux43aux442ux438ux440ux43eux432ux430ux43dux438ux435-ux440ux430ux441ux448ux438ux440ux44fux435ux43cux43eux439-iotml-ux43fux43bux430ux442ux444ux43eux440ux43cux44b-ux443ux43fux440ux430ux432ux43bux435ux43dux438ux44f-ux43cux438ux43aux440ux43eux43aux43bux438ux43cux430ux442ux43eux43c-ux442ux435ux43fux43bux438ux446ux44b}

\subsection{2.1 Назначение проектируемой
платформы}\label{ux43dux430ux437ux43dux430ux447ux435ux43dux438ux435-ux43fux440ux43eux435ux43aux442ux438ux440ux443ux435ux43cux43eux439-ux43fux43bux430ux442ux444ux43eux440ux43cux44b}

Проектируемая система предназначена для мониторинга, журналирования и
управления микроклиматом опытной теплицы. В рамках ВКР она
рассматривается не как частная схема подключения одного датчика к одному
реле, а как расширяемая программно-аппаратная IoT/ML-платформа. Такая
постановка важна, поскольку интеллектуальное управление невозможно без
устойчивого нижнего инженерного слоя: датчики должны регулярно
публиковать данные, исполнительные устройства должны получать команды и
подтверждать состояние, а история измерений и управляющих воздействий
должна сохраняться в форме, пригодной для анализа.

Теплица в работе рассматривается как теплотехнический и киберфизический
объект. Температура, относительная влажность, концентрация CO\(_2\),
вентиляция, нагрев, увлажнение и освещение образуют взаимосвязанную
систему. Изменение одного воздействия может менять сразу несколько
параметров: вентиляция снижает концентрацию CO\(_2\) и влажность, но
может увеличивать теплопотери; нагрев меняет температуру и относительную
влажность; увлажнение влияет на влажностный режим и косвенно на тепловой
баланс. Поэтому архитектура должна сохранять не только значения
датчиков, но и события включения оборудования, иначе последующий анализ
реакции микроклимата будет неполным.

Целевое назначение платформы состоит в том, чтобы обеспечить:

\begin{enumerate}
\def\labelenumi{\arabic{enumi})}
\tightlist
\item
  подключение разнородных сенсорных узлов;
\item
  управление исполнительными устройствами через единый контур команд;
\item
  отображение текущего состояния теплицы оператору;
\item
  хранение истории микроклимата и управляющих воздействий;
\item
  экспорт данных для статистического анализа и машинного обучения;
\item
  возможность развития от ручного и правилового управления к прогнозной
  поддержке решений.
\end{enumerate}

Такой подход согласуется с современными представлениями о smart
greenhouse, где IoT-инфраструктура, web-интерфейс, история данных и
интеллектуальный слой являются частями единой системы
{[}bersani2022iotsmartgreenhouse; elouaham2026smartgreenhouses;
thwin2024microclimateweb{]}.

\subsection{2.2 Требования к
системе}\label{ux442ux440ux435ux431ux43eux432ux430ux43dux438ux44f-ux43a-ux441ux438ux441ux442ux435ux43cux435}

Требования к системе разделены на функциональные, нефункциональные и
исследовательские. Функциональные требования описывают, что система
должна делать в текущей реализации. Нефункциональные требования задают
свойства надежности и сопровождаемости. Исследовательские требования
связаны с тем, что платформа должна формировать данные для анализа и
последующего ML-слоя.

К функциональным требованиям относятся:

\begin{enumerate}
\def\labelenumi{\arabic{enumi})}
\tightlist
\item
  сбор температуры и влажности с сенсорных узлов;
\item
  сбор концентрации CO\(_2\), температуры и влажности с отдельного
  сенсорного узла;
\item
  управление многоканальным релейным контроллером;
\item
  передача телеметрии, команд и статусов через MQTT;
\item
  отображение состояния датчиков, реле и контроллера в web-dashboard;
\item
  настройка правил автоматизации по датчикам, расписанию и гистерезису;
\item
  настройка профилей управления с последовательным выполнением команд и
  задержками;
\item
  хранение истории сенсорных измерений и переключений реле в локальной
  базе данных;
\item
  экспорт истории в форматы, пригодные для анализа.
\end{enumerate}

К нефункциональным требованиям относятся устойчивое переподключение
Wi-Fi и MQTT, публикация статусов online/offline, разделение команд и
подтвержденных состояний реле, сохранение конфигурации в файлах,
простота обслуживания и возможность добавления новых устройств без
полной переработки архитектуры.

Исследовательские требования включают синхронное журналирование
микроклимата и управляющих воздействий, сохранение временных меток в
едином формате, возможность формирования равномерных временных рядов,
построение лаговых признаков и сопоставление включений оборудования с
изменениями температуры, влажности и CO\(_2\). Эти требования
непосредственно связаны с задачей машинного обучения: модель не может
прогнозировать реакцию теплицы, если в данных есть только сенсорные
значения, но нет факта управляющего воздействия.

\subsection{2.3 Общая архитектура
платформы}\label{ux43eux431ux449ux430ux44f-ux430ux440ux445ux438ux442ux435ux43aux442ux443ux440ux430-ux43fux43bux430ux442ux444ux43eux440ux43cux44b}

Архитектура системы построена по распределенному принципу и включает
четыре уровня:

\begin{enumerate}
\def\labelenumi{\arabic{enumi})}
\tightlist
\item
  сенсорный уровень;
\item
  исполнительный уровень;
\item
  коммуникационный уровень;
\item
  серверно-аналитический уровень.
\end{enumerate}

Сенсорный уровень представлен узлами на ESP8266. В текущем инкременте
используются датчик температуры и влажности DHT11 и датчик SCD30,
измеряющий CO\(_2\), температуру и влажность. Исполнительный уровень
представлен контроллером на Arduino Mega 2560 с ESP8266 в режиме AT,
который управляет 15 релейными каналами. Коммуникационный уровень
построен на MQTT-брокере Mosquitto. Серверно-аналитический уровень
реализован web-dashboard, который принимает сообщения, отображает
состояние, публикует команды, выполняет правила, запускает профили и
сохраняет историю в SQLite.

Распределенная структура выбрана по трем причинам. Во-первых, она
позволяет физически разносить датчики по зонам теплицы и постепенно
добавлять новые узлы. Во-вторых, она отделяет опасные исполнительные
действия от пользовательского интерфейса: web-dashboard не переключает
GPIO напрямую, а публикует команду в MQTT, после чего контроллер реле
подтверждает фактическое состояние. В-третьих, MQTT делает поток данных
наблюдаемым: один и тот же топик может использоваться dashboard,
журналом истории, отладочным клиентом и будущим ML-сервисом.

Для финального текста ВКР общую архитектуру целесообразно сопровождать
схемой из файла диаграмм: сенсорные узлы публикуют телеметрию в MQTT,
dashboard подписывается на топики, контроллер реле получает команды, а
SQLite сохраняет историю.

\subsection{2.4 Реестр устройств и
метрик}\label{ux440ux435ux435ux441ux442ux440-ux443ux441ux442ux440ux43eux439ux441ux442ux432-ux438-ux43cux435ux442ux440ux438ux43a}

Для расширяемой платформы важно отделять физическое устройство от
измеряемой метрики. Один узел может публиковать несколько параметров, а
одна и та же метрика может измеряться несколькими устройствами в разных
зонах. Поэтому в проектной модели вводятся две сущности: реестр
устройств и реестр метрик.

Реестр устройств содержит идентификатор, тип устройства, расположение,
набор публикуемых метрик, статус доступности и служебные параметры. В
текущем прототипе используются идентификаторы \texttt{th-1},
\texttt{th-2}, \texttt{co2-1} и \texttt{mega-1}. Узлы \texttt{th-1} и
\texttt{th-2} относятся к датчикам температуры и влажности,
\texttt{co2-1} - к датчику CO\(_2\), температуры и влажности,
\texttt{mega-1} - к контроллеру исполнительных устройств.

Реестр метрик описывает физический смысл поля данных: температура
воздуха, относительная влажность, концентрация CO\(_2\), состояние реле,
длительность включения, признак доступности устройства. Для каждой
метрики должны быть заданы единица измерения, допустимый физический
диапазон и правило обработки некорректных значений. Например,
температура хранится в градусах Цельсия, влажность - в процентах,
CO\(_2\) - в ppm. Нулевые значения влажности и CO\(_2\) в реальных
данных не должны автоматически интерпретироваться как физическое
состояние среды: они могут быть признаком технического сбоя или
отсутствия измерения.

Такое разделение позволяет добавлять новые датчики без изменения смысла
аналитики. Если в систему будет добавлен наружный датчик температуры, он
получит новый идентификатор устройства, но метрика \texttt{temperature}
останется той же. Если будет добавлен датчик освещенности, в реестр
добавится новая метрика, а pipeline машинного обучения сможет
использовать ее как дополнительный признак.

\subsection{2.5 Унифицированная модель
MQTT-обмена}\label{ux443ux43dux438ux444ux438ux446ux438ux440ux43eux432ux430ux43dux43dux430ux44f-ux43cux43eux434ux435ux43bux44c-mqtt-ux43eux431ux43cux435ux43dux430}

В системе используется единое пространство MQTT-топиков с префиксом
\texttt{greenhouse}. Сенсорные узлы публикуют измерения в топики вида:

\begin{verbatim}
greenhouse/env/<device>/state
\end{verbatim}

где \texttt{\textless{}device\textgreater{}} - идентификатор устройства.
Статус доступности публикуется в топики:

\begin{verbatim}
greenhouse/<device>/status
\end{verbatim}

Команды реле передаются через топики:

\begin{verbatim}
greenhouse/relay/<n>/set
\end{verbatim}

Подтвержденные состояния реле публикуются в:

\begin{verbatim}
greenhouse/relay/<n>/state
\end{verbatim}

Отдельно используется сводный топик \texttt{greenhouse/relay/summary}, в
котором контроллер передает диагностические сведения: время работы,
состояние Wi-Fi и MQTT, уровень сигнала и массив состояний реле.

Сенсорное сообщение передается в JSON-формате. Для текущих устройств
фактически используются поля \texttt{ts}, \texttt{device}, \texttt{t},
\texttt{h} и \texttt{co2}. В платформенной модели этот формат можно
обобщить следующим образом:

\begin{verbatim}
{
  "ts": 123456,
  "device": "co2-1",
  "metrics": {
    "temperature": 24.7,
    "humidity": 60.5,
    "co2": 820
  }
}
\end{verbatim}

Текущий прототип использует более компактный формат для
микроконтроллеров, но проектная логика остается такой же: каждое
сообщение содержит источник, время и набор метрик. Для совместимости
dashboard может поддерживать оба варианта: короткие поля для
существующих прошивок и обобщенный словарь \texttt{metrics} для новых
устройств.

Главное архитектурное требование к обмену состоит в разделении команды и
фактического состояния. Команда \texttt{ON} в топике \texttt{set}
означает намерение включить реле, но истинным состоянием считается
только сообщение контроллера в топике \texttt{state}. Это снижает риск
рассинхронизации интерфейса и оборудования.

\subsection{2.6 Нижний контур
управления}\label{ux43dux438ux436ux43dux438ux439-ux43aux43eux43dux442ux443ux440-ux443ux43fux440ux430ux432ux43bux435ux43dux438ux44f}

Нижний контур управления включает сенсорные узлы, контроллер реле и
правила, выполняемые серверным приложением. Он должен оставаться
простым, предсказуемым и безопасным. Поэтому текущая реализация
использует ручные команды, пороговые правила, расписания, гистерезис и
профили последовательного включения оборудования.

Сенсорные узлы выполняют циклическое чтение датчиков, проверку
корректности значений, подключение к Wi-Fi, подключение к MQTT и
публикацию статуса доступности. Для CO\(_2\) предусмотрена фильтрация
явно некорректного диапазона. Это не заменяет последующую очистку
данных, но уменьшает попадание грубых выбросов в поток телеметрии.

Контроллер реле принимает команды \texttt{ON}, \texttt{OFF} и
\texttt{TOGGLE}, меняет состояние физического канала и публикует
подтверждение. Он также публикует диагностическую сводку, которая нужна
оператору и журналу состояния. В текущей конфигурации 15 релейных
каналов могут использоваться для насосов, фильтров, вентиляторов,
заслонок, нагревателей и увлажнителей.

Правила и профили находятся выше микроконтроллеров, но ниже будущего
ML-слоя. Правило проверяет значение датчика, расписание и гистерезис.
Профиль задает последовательность действий, например сначала открыть
заслонки, затем после задержки включить приток или вентиляцию. При
остановке профиля последовательность выполняется в обратном порядке:
сначала выключаются основные устройства, затем закрываются
предварительно открытые каналы. Такая логика отражает физическую
последовательность работы оборудования и снижает риск некорректного
порядка включения.

\subsection{2.7 Серверный слой и хранение
истории}\label{ux441ux435ux440ux432ux435ux440ux43dux44bux439-ux441ux43bux43eux439-ux438-ux445ux440ux430ux43dux435ux43dux438ux435-ux438ux441ux442ux43eux440ux438ux438}

Серверный слой выполняет функции интерфейса оператора, MQTT-клиента,
исполнителя правил, менеджера профилей и регистратора истории. Dashboard
отображает последние значения датчиков, состояния реле, статус
контроллера и доступность устройств. При ручном управлении оператор
инициирует команду через web-интерфейс, но физическое переключение
происходит только после получения команды контроллером реле.

История хранится в SQLite. В текущей реализации используются две
ключевые таблицы: \texttt{sensor\_log} и \texttt{relay\_log}. Таблица
\texttt{sensor\_log} содержит временную метку, устройство, температуру,
влажность и CO\(_2\). Таблица \texttt{relay\_log} содержит временную
метку, номер реле и состояние после переключения. Такая минимальная
структура уже достаточна для формирования временных рядов и анализа
естественных экспериментов.

Для развития платформы целесообразно расширить модель данных до
следующих сущностей:

{\def\LTcaptype{none} % do not increment counter
\begin{longtable}[]{@{}
  >{\RaggedRight\arraybackslash}p{(\linewidth - 2\tabcolsep) * \real{0.5000}}
  >{\RaggedRight\arraybackslash}p{(\linewidth - 2\tabcolsep) * \real{0.5000}}@{}}
\hline
\begin{minipage}[b]{\linewidth}\RaggedRight
Сущность
\end{minipage} & \begin{minipage}[b]{\linewidth}\RaggedRight
Назначение
\end{minipage} \\
\hline
\endhead
\hline
\endlastfoot
\texttt{devices} & реестр сенсорных и исполнительных устройств \\
\texttt{metrics} & справочник измеряемых параметров и единиц \\
\texttt{sensor\_readings} & нормализованная история измерений \\
\texttt{actuator\_events} & события исполнительных устройств \\
\texttt{control\_profiles} & профили управления и последовательности
действий \\
\texttt{experiment\_batches} & партии растений и сроки выращивания \\
\texttt{growth\_variants} & варианты света, питания, зоны и ярусы \\
\texttt{plant\_measurements} & биометрия, хлорофилл, качество и
хранение \\
\texttt{ml\_feature\_windows} & подготовленные окна признаков для
обучения \\
\end{longtable}
}

Текущие таблицы можно рассматривать как первый инкремент этой модели.
Они не исчерпывают будущую платформу, но уже фиксируют главное:
совместную временную шкалу микроклимата и управляющих воздействий.

\subsection{2.8 Модель данных эксперимента
выращивания}\label{ux43cux43eux434ux435ux43bux44c-ux434ux430ux43dux43dux44bux445-ux44dux43aux441ux43fux435ux440ux438ux43cux435ux43dux442ux430-ux432ux44bux440ux430ux449ux438ux432ux430ux43dux438ux44f}

Поскольку теплица используется не только как технический объект, но и
как среда выращивания растений, платформа должна поддерживать связь
микроклимата с агробиологическими данными. Для салата важны фотопериод,
PPFD, DLI, электрическая мощность освещения, EC и уровень питательного
раствора, температура раствора, биометрические измерения, содержание
хлорофилла, признаки стеблевания, потери массы при хранении и
органолептическая оценка.

В проектной модели партия растений задается датой посева или высадки,
сортом, вариантом выращивания, зоной размещения и световым режимом.
Вариант выращивания описывает фотопериод, PPFD, DLI и электрическую
мощность светильника. Агробиологические измерения привязываются к
партии, дате и варианту. Микроклиматические данные привязываются не к
отдельному растению, а к временному окну и зоне, в которой находилась
партия.

Такое разделение важно для корректной интерпретации. Нельзя напрямую
утверждать, что разница массы растений вызвана только температурой или
только светом, если факторы не были специально рандомизированы и
синхронизированы. Но платформа позволяет подготовить следующую
постановку controlled experiment: для каждой партии можно заранее
зафиксировать световой вариант, зону, период выращивания, набор датчиков
и временные окна, по которым будут агрегироваться температура,
влажность, CO\(_2\) и состояния исполнительных устройств.

\subsection{2.9 Математический каркас тепловлажностного
режима}\label{ux43cux430ux442ux435ux43cux430ux442ux438ux447ux435ux441ux43aux438ux439-ux43aux430ux440ux43aux430ux441-ux442ux435ux43fux43bux43eux432ux43bux430ux436ux43dux43eux441ux442ux43dux43eux433ux43e-ux440ux435ux436ux438ux43cux430}

Проектируемая платформа не заменяет полную физическую модель теплицы,
однако для магистерской работы необходим формальный язык описания
микроклимата. В дальнейшем этот язык используется для интерпретации
данных и постановки ML-задачи.

Упрощенный тепловой баланс теплицы можно записать в виде:

\[
C_T \frac{dT}{dt} = Q_{heat} + Q_{light} + Q_{solar} - Q_{loss} - Q_{vent} - Q_{evap},
\]

где \texttt{T} - температура воздуха, \texttt{C\_T} - эффективная
теплоемкость воздуха и конструкций, \texttt{Q\_heat} - тепловая мощность
нагрева, \texttt{Q\_light} - тепловыделение освещения, \texttt{Q\_solar}
- приток солнечной радиации, \texttt{Q\_loss} - теплопотери через
ограждения, \texttt{Q\_vent} - потери или приток тепла при вентиляции,
\texttt{Q\_evap} - тепловой эффект испарения.

Баланс влаги можно представить как:

\[
C_H \frac{dH}{dt} = G_{evap} + G_{humid} - G_{vent} - G_{cond},
\]

где \texttt{H} - влажностная характеристика воздуха, \texttt{G\_evap} -
испарение с растений и раствора, \texttt{G\_humid} - подача влаги
увлажнителем, \texttt{G\_vent} - удаление влаги вентиляцией,
\texttt{G\_cond} - конденсация.

Баланс CO\(_2\) в упрощенном виде:

\[
V \frac{dC}{dt} = F_{in}(C_{out} - C) + G_{CO$_2$} - U_{plant},
\]

где \texttt{C} - концентрация CO\(_2\) в теплице, \texttt{V} - объем
воздуха, \texttt{F\_in} - воздухообмен, \texttt{C\_out} - наружная
концентрация, \texttt{G\_CO\$\_2\$} - подача или внутреннее образование
CO\(_2\), \texttt{U\_plant} - поглощение растениями.

Эти уравнения показывают, почему в данных нужны не только значения
датчиков, но и управляющие воздействия. Без информации о вентиляции,
нагреве и увлажнении невозможно отличить естественный суточный ход от
реакции на оборудование.

Для анализа временных рядов используются более простые расчетные
показатели. Изменение параметра после события реле определяется как:

\[
\Delta x = \bar{x}_{after} - \bar{x}_{before},
\]

где \texttt{\textbackslash{}bar\{x\}\_\{before\}} - среднее значение в
окне до события, а \texttt{\textbackslash{}bar\{x\}\_\{after\}} -
среднее значение в окне после события. Для прогноза используются
стандартные метрики:

\[
MAE = \frac{1}{n}\sum_{i=1}^{n}|y_i-\hat{y}_i|,
\]

\[
RMSE = \sqrt{\frac{1}{n}\sum_{i=1}^{n}(y_i-\hat{y}_i)^2},
\]

\[
R^2 = 1 - \frac{\sum_{i=1}^{n}(y_i-\hat{y}_i)^2}{\sum_{i=1}^{n}(y_i-\bar{y})^2}.
\]

\subsection{2.10 Подготовка данных для
ML-слоя}\label{ux43fux43eux434ux433ux43eux442ux43eux432ux43aux430-ux434ux430ux43dux43dux44bux445-ux434ux43bux44f-ml-ux441ux43bux43eux44f}

ML-слой в работе рассматривается как верхний прогнозный уровень, а не
как замена безопасного правилового управления. Его первая задача -
краткосрочный прогноз температуры, влажности и CO\(_2\) на горизонте
10-30 минут. Такой прогноз может использоваться для раннего
предупреждения о перегреве, переувлажнении или снижении CO\(_2\), а в
дальнейшем - для рекомендации уставок и выбора профиля.

Подготовка ML-таблицы включает несколько шагов:

\begin{enumerate}
\def\labelenumi{\arabic{enumi})}
\tightlist
\item
  выгрузку \texttt{sensor\_log} и \texttt{relay\_log};
\item
  приведение временных меток к единой шкале;
\item
  разворот сенсорных данных в wide-формат по устройствам и метрикам;
\item
  ресемплирование к шагу 1-5 минут;
\item
  восстановление состояний реле на временной шкале;
\item
  фильтрацию технических нулей и физических выбросов;
\item
  построение лагов 5, 10, 30 и 60 минут;
\item
  расчет скользящих средних и стандартных отклонений;
\item
  добавление календарных признаков и признаков день/ночь;
\item
  формирование целевых переменных на выбранном горизонте.
\end{enumerate}

Для первой проверки достаточно сравнить naive baseline, линейную
регрессию и табличную ML-модель. Более сложные LSTM, GRU, TCN или
Transformer целесообразны после расширения датасета и добавления внешних
факторов, поскольку нейросетевые модели требуют более строгой валидации
и большего разнообразия режимов {[}ahn2024transformerrnngreenhouse;
gong2021tcnrnnyield; choi2025probabilisticmicroclimate{]}.

\subsection{2.11 Расчетная модель экономического
эффекта}\label{ux440ux430ux441ux447ux435ux442ux43dux430ux44f-ux43cux43eux434ux435ux43bux44c-ux44dux43aux43eux43dux43eux43cux438ux447ux435ux441ux43aux43eux433ux43e-ux44dux444ux444ux435ux43aux442ux430}

Экономический эффект в работе формулируется как расчетный потенциальный
эффект, а не как экспериментально доказанный прирост прибыли. Это
принципиальное ограничение: имеющийся датасет показывает работу
платформы и реакцию микроклимата, но не является полнофакторным
экономическим экспериментом.

Расчетная модель включает несколько составляющих:

\begin{enumerate}
\def\labelenumi{\arabic{enumi})}
\tightlist
\item
  энергозатраты освещения;
\item
  стоимость световой нагрузки DLI;
\item
  энергозатраты исполнительных устройств по журналу реле;
\item
  сценарную экономию от прогнозного управления;
\item
  потенциальное снижение потерь товарной продукции;
\item
  экономию трудозатрат оператора;
\item
  срок окупаемости оборудования и внедрения.
\end{enumerate}

Энергозатраты освещения определяются по формуле:

\[
E_{light} = P_{light} \cdot t_{light},
\]

где \texttt{P\_light} - электрическая мощность светильника,
\texttt{t\_light} - продолжительность работы. Для световых вариантов с
фотопериодом 18 ч/сут DLI рассчитывается как:

\[
DLI = \frac{PPFD \cdot 18 \cdot 3600}{1000000}.
\]

Энергия исполнительных устройств рассчитывается по длительности
включения реле:

\[
E_{act} = \sum_{j=1}^{m} P_j \cdot t_j,
\]

где \texttt{P\_j} - номинальная мощность устройства, \texttt{t\_j} -
суммарное время включения по журналу реле. Если фактическая мощность
части оборудования неизвестна, расчет должен выполняться сценарно с явно
указанными допущениями.

Сценарная экономия от прогнозного управления может быть оценена как доля
предотвращенных лишних включений:

\[
S_{energy} = E_{base} \cdot k_{save} \cdot c_{el},
\]

где \texttt{E\_base} - базовое энергопотребление, \texttt{k\_save} -
сценарная доля экономии, \texttt{c\_el} - тариф на электроэнергию.
Отдельно может оцениваться эффект от снижения потерь продукции и
сокращения ручного контроля, но эти величины должны быть представлены
как расчетная оценка при заданных предпосылках.

\subsection{2.12 Выводы по
главе}\label{ux432ux44bux432ux43eux434ux44b-ux43fux43e-ux433ux43bux430ux432ux435}

Во второй главе спроектирована расширяемая IoT/ML-платформа мониторинга
и управления микроклиматом теплицы. В отличие от простой схемы локальной
автоматики, архитектура включает сенсорный слой, исполнительный
контроллер, MQTT-обмен, web-dashboard, правила, профили, журналирование
истории и подготовку данных для машинного обучения.

Ключевой проектный результат состоит в том, что платформа разделяет
устройства, метрики, команды, подтвержденные состояния и исторические
события. Благодаря этому новые датчики и исполнительные каналы могут
подключаться через общую модель обмена, а накопленная история становится
пригодной для анализа реакции теплицы, расчета естественных
экспериментов и построения прогнозных моделей.

Предложенный математический и экономический каркас задает язык
дальнейшей практической главы: фактический датасет будет описываться как
временной ряд микроклимата и управляющих воздействий, реакции на
включения оборудования - как наблюдательные естественные эксперименты, а
экономический эффект - как сценарная расчетная оценка потенциальной
пользы прогнозного управления.

\section{3 Реализация, опытное внедрение и анализ данных
платформы}\label{ux440ux435ux430ux43bux438ux437ux430ux446ux438ux44f-ux43eux43fux44bux442ux43dux43eux435-ux432ux43dux435ux434ux440ux435ux43dux438ux435-ux438-ux430ux43dux430ux43bux438ux437-ux434ux430ux43dux43dux44bux445-ux43fux43bux430ux442ux444ux43eux440ux43cux44b}

\subsection{3.1 Состав реализованной
системы}\label{ux441ux43eux441ux442ux430ux432-ux440ux435ux430ux43bux438ux437ux43eux432ux430ux43dux43dux43eux439-ux441ux438ux441ux442ux435ux43cux44b}

Практическая часть ВКР реализована как распределенная IoT-система для
опытной теплицы. В нее входят сенсорные узлы, контроллер исполнительных
устройств, MQTT-брокер, web-dashboard, правила автоматизации, профили
управления и локальная база истории. Такой состав соответствует
проектной архитектуре, описанной во второй главе, и позволяет
рассматривать систему как первый инкремент расширяемой IoT/ML-платформы.

Фактическая реализация включает:

\begin{enumerate}
\def\labelenumi{\arabic{enumi})}
\tightlist
\item
  сенсорный узел температуры и влажности на ESP8266 и DHT11;
\item
  сенсорный узел CO\(_2\), температуры и влажности на ESP8266 и SCD30;
\item
  релейный контроллер на Arduino Mega 2560 с ESP8266 AT;
\item
  MQTT-брокер Mosquitto;
\item
  Flask-dashboard оператора;
\item
  SQLite-базу истории;
\item
  JSON-конфигурации правил, профилей и имен исполнительных каналов.
\end{enumerate}

Основная инженерная особенность реализации состоит в разделении функций.
Сенсорные узлы не принимают решений и не управляют реле напрямую.
Релейный контроллер не вычисляет правила микроклимата, а исполняет
команды и подтверждает состояние. Серверный слой объединяет телеметрию,
интерфейс оператора, правила, профили и журналирование. Такое разделение
снижает связанность компонентов и упрощает дальнейшее расширение
платформы.

\subsection{3.2 Реализация сенсорных
узлов}\label{ux440ux435ux430ux43bux438ux437ux430ux446ux438ux44f-ux441ux435ux43dux441ux43eux440ux43dux44bux445-ux443ux437ux43bux43eux432}

Узел температуры и влажности использует ESP8266 и датчик DHT11. Он
подключается к Wi-Fi, устанавливает MQTT-соединение, публикует статус
доступности и отправляет JSON-сообщения в топик
\texttt{greenhouse/env/th-1/state}. Сообщение содержит временную метку,
идентификатор устройства, температуру и влажность:

\begin{verbatim}
{"ts":123456,"device":"th-1","t":24.5,"h":61.0}
\end{verbatim}

Если датчик возвращает некорректное значение, публикация пропускается.
Это важно для дальнейшего анализа, поскольку нечисловые значения не
должны попадать в историю как физические измерения.

Узел CO\(_2\) использует ESP8266 и датчик SCD30 по I2C. Он публикует
данные в \texttt{greenhouse/env/co2-1/state} и передает концентрацию
CO\(_2\), температуру и влажность:

\begin{verbatim}
{"ts":123456,"device":"co2-1","co2":820,"t":24.7,"h":60.5}
\end{verbatim}

В прошивке предусмотрена проверка диапазона CO\(_2\). Значения меньше 0
ppm и больше 10000 ppm отбрасываются, после чего выполняется попытка
восстановления работы датчика. Это не устраняет все проблемы качества
данных, но уменьшает число грубых выбросов на входе платформы.

В опытном датасете также присутствует устройство \texttt{th-2}. Его
наличие показывает, что архитектура уже работает не как жесткая схема
одного датчика, а как модель с несколькими источниками измерений,
которые можно сравнивать между собой.

\subsection{3.3 Реализация контроллера
реле}\label{ux440ux435ux430ux43bux438ux437ux430ux446ux438ux44f-ux43aux43eux43dux442ux440ux43eux43bux43bux435ux440ux430-ux440ux435ux43bux435}

Исполнительный уровень реализован на Arduino Mega 2560 с ESP8266 в
режиме AT. Контроллер управляет 15 релейными каналами. Команды поступают
через MQTT-топики
\texttt{greenhouse/relay/\textless{}n\textgreater{}/set}, где
\texttt{\textless{}n\textgreater{}} - номер реле. Поддерживаются команды
\texttt{ON}, \texttt{OFF}, \texttt{TOGGLE}, а также числовые варианты
\texttt{1} и \texttt{0}. После изменения состояния контроллер публикует
подтверждение в
\texttt{greenhouse/relay/\textless{}n\textgreater{}/state}.

Такой механизм принципиален для надежного управления. Dashboard не
считает реле включенным только потому, что оператор нажал кнопку.
Истинным состоянием считается подтверждение от контроллера. Это снижает
риск ошибочного отображения состояния вентилятора, заслонки, увлажнителя
или нагревателя.

Контроллер также публикует диагностический топик
\texttt{greenhouse/relay/summary}, в котором передаются состояние связи,
время работы и массив состояний реле. При подключении используется
статус \texttt{online}, а при аварийном разрыве MQTT-брокер может
отметить устройство как \texttt{offline}. Эти статусы нужны не только
оператору, но и анализу качества данных: отсутствие показаний или
подтверждений может быть связано с сетью, а не с устойчивым состоянием
микроклимата.

В текущей конфигурации наиболее важными для анализа оказались реле 8, 9
и 11, связанные с профилем притока с улицы для бокса 2. Также в журнале
есть события реле 2 и 3, относящиеся к ультразвуковому увлажнению, и
отдельные короткие события других каналов.

\subsection{3.4 Dashboard, правила и профили
управления}\label{dashboard-ux43fux440ux430ux432ux438ux43bux430-ux438-ux43fux440ux43eux444ux438ux43bux438-ux443ux43fux440ux430ux432ux43bux435ux43dux438ux44f}

Серверная часть реализована как Flask-dashboard. Приложение
подписывается на MQTT-топики датчиков, реле и статусов, хранит последние
значения в оперативном состоянии, отображает их в web-интерфейсе и
публикует команды управления. Пользовательский интерфейс содержит
главную страницу состояния, страницу правил, страницу профилей и
страницу статистики.

Правила автоматизации задают условие по датчику, расписание, гистерезис
и действие. Например, правило может включить профиль вентиляции при
превышении порога температуры или CO\(_2\) в заданный временной
интервал. Гистерезис нужен для предотвращения частых переключений около
порога.

Профили управления задают последовательность действий. Для вентиляции
типовой порядок таков: сначала открыть заслонки, затем после задержки
включить вентиляторы или приток. При выключении профиль выполняется в
обратной логике: сначала отключаются основные устройства, затем
закрываются предварительно открытые каналы. Такой порядок отражает
физику оборудования и делает правила безопаснее, чем простое
одновременное включение всех реле.

С точки зрения интеллектуального управления правила и профили являются
нижним уровнем. Они остаются интерпретируемыми и пригодными для ручной
проверки. ML-слой в дальнейшем должен не заменять их напрямую, а
прогнозировать риск выхода параметров за допустимые границы и помогать
выбирать момент запуска профиля.

\subsection{3.5 Журналирование и экспорт
истории}\label{ux436ux443ux440ux43dux430ux43bux438ux440ux43eux432ux430ux43dux438ux435-ux438-ux44dux43aux441ux43fux43eux440ux442-ux438ux441ux442ux43eux440ux438ux438}

Для хранения истории используется SQLite. Таблица \texttt{sensor\_log}
хранит временную метку, устройство, температуру, влажность и CO\(_2\).
Таблица \texttt{relay\_log} хранит временную метку, номер реле и
состояние после события. Сенсорные значения записываются с
периодичностью, достаточной для анализа микроклимата, а события реле
фиксируются при изменении состояния.

По результатам опытной эксплуатации база была скопирована с Raspberry Pi
\texttt{green} и сохранена как локальное доказательство внедрения.
Диапазон данных:

{\def\LTcaptype{none} % do not increment counter
\begin{longtable}[]{@{}lrl@{}}
\hline
Таблица & Число строк & Период \\
\hline
\endhead
\hline
\endlastfoot
\texttt{sensor\_log} & 243946 & 28.03.2026 16:12:15 - 26.05.2026
22:26:46 \\
\texttt{relay\_log} & 165 & 28.03.2026 16:22:35 - 25.05.2026 11:30:07 \\
\end{longtable}
}

Датасет охватывает почти два месяца опытной эксплуатации. Это уже
достаточная основа для описательного анализа, проверки качества данных,
выделения естественных экспериментов и постановки baseline-прогноза. При
этом для строгого факторного вывода данных пока недостаточно, поскольку
отсутствуют наружная погода, точные режимы всех внешних факторов и
специально рандомизированный план воздействий.

\subsection{3.6 Фактический датасет
микроклимата}\label{ux444ux430ux43aux442ux438ux447ux435ux441ux43aux438ux439-ux434ux430ux442ux430ux441ux435ux442-ux43cux438ux43aux440ux43eux43aux43bux438ux43cux430ux442ux430}

В датасете представлены три сенсорных источника: \texttt{th-1},
\texttt{th-2} и \texttt{co2-1}. Сводные характеристики после исключения
технических нулей показывают, что температура по датчикам находилась в
диапазоне примерно 17,6-34,9 °C, а средние значения различались по
точкам измерения.

{\def\LTcaptype{none} % do not increment counter
\begin{longtable}[]{@{}
  >{\RaggedRight\arraybackslash}p{(\linewidth - 8\tabcolsep) * \real{0.1579}}
  >{\Centering\arraybackslash}p{(\linewidth - 8\tabcolsep) * \real{0.2105}}
  >{\Centering\arraybackslash}p{(\linewidth - 8\tabcolsep) * \real{0.2105}}
  >{\Centering\arraybackslash}p{(\linewidth - 8\tabcolsep) * \real{0.2105}}
  >{\Centering\arraybackslash}p{(\linewidth - 8\tabcolsep) * \real{0.2105}}@{}}
\hline
\begin{minipage}[b]{\linewidth}\RaggedRight
Устройство
\end{minipage} & \begin{minipage}[b]{\linewidth}\Centering
Строк
\end{minipage} & \begin{minipage}[b]{\linewidth}\Centering
Средняя температура, °C
\end{minipage} & \begin{minipage}[b]{\linewidth}\Centering
Средняя влажность, \%
\end{minipage} & \begin{minipage}[b]{\linewidth}\Centering
Средний CO\(_2\), ppm
\end{minipage} \\
\hline
\endhead
\hline
\endlastfoot
\texttt{co2-1} & 81315 & 24,98 & 37,36 & 354,61 \\
\texttt{th-1} & 81316 & 25,17 & 28,53 & - \\
\texttt{th-2} & 81315 & 27,85 & 27,57 & - \\
\end{longtable}
}

Различия между датчиками имеют практическое значение. Среднее абсолютное
расхождение температуры между \texttt{th-1} и \texttt{th-2} составило
около 2,69 °C, а между \texttt{th-2} и \texttt{co2-1} - около 3,23 °C.
Это может отражать как пространственную неоднородность микроклимата, так
и смещение датчиков. Для платформенной ВКР такой результат важен: один
датчик не описывает всю теплицу, поэтому архитектура должна поддерживать
несколько точек измерения и хранить идентификатор источника.

Анализ качества данных выявил технические нули у \texttt{co2-1}: 28752
строки с нулевой температурой, 28752 строки с нулевой влажностью и 28753
строки с нулевым CO\(_2\). Также у \texttt{th-2} найдено 36 нулевых
значений влажности. Эти значения не следует использовать как физические
измерения. Они сохранены как признак качества данных и должны
исключаться или маркироваться при расчетах и обучении моделей.

В данных есть крупный разрыв около 3965 минут в начале периода: с
28.03.2026 примерно 17:26 до 31.03.2026 примерно 11:31. Этот факт
ограничивает анализ ранних событий реле и требует осторожности при
расчете интервалов до/после.

\subsection{3.7 Журнал реле и естественные
эксперименты}\label{ux436ux443ux440ux43dux430ux43b-ux440ux435ux43bux435-ux438-ux435ux441ux442ux435ux441ux442ux432ux435ux43dux43dux44bux435-ux44dux43aux441ux43fux435ux440ux438ux43cux435ux43dux442ux44b}

Журнал реле содержит 165 событий. Наиболее активными оказались каналы 8,
9 и 11, связанные с профилем притока с улицы для бокса 2. Их активность
делает возможным наблюдательный анализ реакции микроклимата на включения
оборудования.

{\def\LTcaptype{none} % do not increment counter
\begin{longtable}[]{@{}
  >{\Centering\arraybackslash}p{(\linewidth - 10\tabcolsep) * \real{0.1739}}
  >{\RaggedRight\arraybackslash}p{(\linewidth - 10\tabcolsep) * \real{0.1304}}
  >{\Centering\arraybackslash}p{(\linewidth - 10\tabcolsep) * \real{0.1739}}
  >{\Centering\arraybackslash}p{(\linewidth - 10\tabcolsep) * \real{0.1739}}
  >{\Centering\arraybackslash}p{(\linewidth - 10\tabcolsep) * \real{0.1739}}
  >{\Centering\arraybackslash}p{(\linewidth - 10\tabcolsep) * \real{0.1739}}@{}}
\hline
\begin{minipage}[b]{\linewidth}\Centering
Реле
\end{minipage} & \begin{minipage}[b]{\linewidth}\RaggedRight
Интерпретация по профилям
\end{minipage} & \begin{minipage}[b]{\linewidth}\Centering
Событий
\end{minipage} & \begin{minipage}[b]{\linewidth}\Centering
ON
\end{minipage} & \begin{minipage}[b]{\linewidth}\Centering
OFF
\end{minipage} & \begin{minipage}[b]{\linewidth}\Centering
Наблюдаемое время ON, мин
\end{minipage} \\
\hline
\endhead
\hline
\endlastfoot
8 & приток с улицы, бокс 2, основное действие & 43 & 22 & 21 &
39556,75 \\
9 & приток с улицы, бокс 2, предварительное действие & 37 & 19 & 18 &
38273,87 \\
11 & приток с улицы, бокс 2, предварительное действие & 33 & 17 & 16 &
39513,27 \\
2 & 2-й ультразвуковой увлажнитель & 18 & 9 & 9 & 2961,92 \\
3 & 1-й ультразвуковой увлажнитель & 6 & 3 & 3 & 1332,62 \\
\end{longtable}
}

Для анализа естественных экспериментов каждое включение реле
рассматривается как событие, вокруг которого строятся окна до и после. В
текущем отчете использовались окна 30 минут до события, 0-30 минут после
события и 30-120 минут после события. Для каждого окна рассчитывались
средние значения температуры, влажности и CO\(_2\), после чего
определялась разница \texttt{после\ -\ до}.

Пример раннего события 28.03.2026 16:54 показывает, что после включения
каналов, связанных с вентиляцией бокса 1, температура по \texttt{th-1} в
первые 30 минут снизилась примерно на 0,51 °C, а через 30-120 минут -
примерно на 1,65 °C. Влажность по \texttt{th-1} в тот же период
увеличилась примерно на 2,31 процентного пункта в первые 30 минут и на
7,0 пункта через 30-120 минут. CO\(_2\) по \texttt{co2-1} снизился
примерно на 54 ppm в первые 30 минут.

Эти результаты нельзя трактовать как строгий причинный эффект, потому
что события могли совпадать с суточным ходом, изменением внешних
условий, несколькими одновременными реле и начальными разрывами данных.
Однако они показывают главное: платформа фиксирует управляющие
воздействия и микроклимат на общей временной шкале, поэтому позволяет
выделять естественные эксперименты и формировать гипотезы для следующего
controlled experiment.

\subsection{3.8 Суточная динамика и лаговые
связи}\label{ux441ux443ux442ux43eux447ux43dux430ux44f-ux434ux438ux43dux430ux43cux438ux43aux430-ux438-ux43bux430ux433ux43eux432ux44bux435-ux441ux432ux44fux437ux438}

Анализ среднесуточных графиков температуры, влажности и CO\(_2\)
показывает выраженную динамику микроклимата. Для теплицы это ожидаемо:
температура и влажность зависят от освещения, воздухообмена, испарения,
нагрева и внешних условий. С практической точки зрения суточная динамика
важна тем, что простое сравнение ``до/после'' без учета времени суток
может дать ложную интерпретацию.

Лаговые корреляции использовались как предварительный инструмент для
поиска связей между рядами. Например, корреляция между CO\(_2\) и
температурой \texttt{co2-1} в разных сдвигах находилась около 0,27-0,28
на ранних проверенных лагах. Это не является сильной связью, но
показывает, что параметры нельзя рассматривать полностью независимо. Для
ML-модели это означает необходимость использовать лаговые признаки и
признаки времени суток.

В дальнейшей версии анализа целесообразно отдельно выделить дневные и
ночные интервалы, поскольку фотосинтез, освещение, вентиляция и
испарение в них имеют разный физический смысл. Также необходимо добавить
наружную температуру и влажность, иначе влияние вентиляции будет трудно
отделить от внешнего климата.

\subsection{3.9 Связь с данными выращивания
салата}\label{ux441ux432ux44fux437ux44c-ux441-ux434ux430ux43dux43dux44bux43cux438-ux432ux44bux440ux430ux449ux438ux432ux430ux43dux438ux44f-ux441ux430ux43bux430ux442ux430}

Отдельным источником данных является таблица по выращиванию салата. Она
содержит сведения о световых вариантах, электрической мощности, EC и
уровне питательного раствора, ручных измерениях температуры раствора,
биометрии, хранении и качественных наблюдениях.

Для световых вариантов зафиксирован фотопериод 18 часов света и 6 часов
темноты. PPFD и рассчитанный DLI составляют:

{\def\LTcaptype{none} % do not increment counter
\begin{longtable}[]{@{}
  >{\Centering\arraybackslash}p{(\linewidth - 6\tabcolsep) * \real{0.2500}}
  >{\Centering\arraybackslash}p{(\linewidth - 6\tabcolsep) * \real{0.2500}}
  >{\Centering\arraybackslash}p{(\linewidth - 6\tabcolsep) * \real{0.2500}}
  >{\Centering\arraybackslash}p{(\linewidth - 6\tabcolsep) * \real{0.2500}}@{}}
\hline
\begin{minipage}[b]{\linewidth}\Centering
Вариант
\end{minipage} & \begin{minipage}[b]{\linewidth}\Centering
PPFD, мкмоль/(м2·с)
\end{minipage} & \begin{minipage}[b]{\linewidth}\Centering
DLI, моль/(м2·сут)
\end{minipage} & \begin{minipage}[b]{\linewidth}\Centering
Электрическая мощность, Вт
\end{minipage} \\
\hline
\endhead
\hline
\endlastfoot
1 & 79,0 & 5,12 & 116,3 \\
2 & 95,0 & 6,16 & 97,0 \\
3 & 150,8 & 9,77 & 221,0 \\
4 & 181,3 & 11,75 & 321,5 \\
5 & 245,5 & 15,91 & 478,0 \\
\end{longtable}
}

Важно не смешивать PPFD и электрическую мощность. PPFD описывает
фотосинтетический поток на уровне растений, а мощность в ваттах -
энергетическую нагрузку оборудования. DLI связывает световую нагрузку с
фотопериодом и рассчитывается по формуле, приведенной во второй главе.

Данные по хранению показывают, что потери массы различаются по сортам и
вариантам. Например, для сорта ``Триплекс'' потери массы составляли
примерно от 12,61 до 21,62 \%, тогда как для ``Джипси'' в видимых
строках таблицы они находились примерно в диапазоне 3,49-4,91 \%. Эти
сведения полезны для экономического блока, но текущая связь с
микроклиматом ограничена: нет точной непрерывной привязки каждого
растения к зоне микроклимата и времени всех ручных измерений.

Следовательно, данные по салату в текущей ВКР следует использовать как
основу экспериментальной рамки, а не как доказательство оптимального
режима выращивания. Платформа показывает, какие данные нужно синхронно
собирать в следующем цикле: световой вариант, зона, партия, микроклимат,
EC, температура раствора, биометрия и итоговое качество.

\subsection{3.10 Расчетный экономический
блок}\label{ux440ux430ux441ux447ux435ux442ux43dux44bux439-ux44dux43aux43eux43dux43eux43cux438ux447ux435ux441ux43aux438ux439-ux431ux43bux43eux43a}

Экономический эффект оценивается расчетно и сценарно. В текущих данных
есть две надежные опоры: электрическая мощность световых вариантов и
журнал длительности включения реле. Этого достаточно, чтобы показать
методику расчета энергозатрат, но недостаточно, чтобы заявлять
доказанную фактическую экономию.

Для освещения суточное потребление при фотопериоде 18 часов можно
оценить как произведение мощности на время. Например, для варианта 5 с
мощностью 478 Вт суточное потребление составляет:

\[
E = 0{,}478 \cdot 18 = 8{,}604 \text{ кВт·ч/сут}.
\]

Для варианта 2 с мощностью 97 Вт:

\[
E = 0{,}097 \cdot 18 = 1{,}746 \text{ кВт·ч/сут}.
\]

При этом вариант 5 дает DLI около 15,91 моль/(м2·сут), а вариант 2 -
около 6,16 моль/(м2·сут). Поэтому экономическое сравнение должно
учитывать не только потребленную энергию, но и световую нагрузку, рост и
качество продукции.

Для исполнительных устройств расчет строится по журналу реле:

\[
E_{relay} = \sum P_j \cdot t_j.
\]

Поскольку фактические мощности части реле требуют уточнения, в ВКР
корректно использовать сценарии. Например, если приточный вентилятор или
группа вентиляции имеет мощность \texttt{P}, то экономия от уменьшения
лишней работы на 10 \% оценивается как \texttt{0,1\ *\ P\ *\ t\ *\ c},
где \texttt{t} - наблюдаемая длительность включения, а \texttt{c} -
тариф. Такой расчет показывает потенциальную ценность прогнозного
управления, но не подменяет собой экспериментальное подтверждение.

Дополнительные составляющие эффекта связаны с потерями товарной массы и
трудозатратами оператора. Если платформа снижает число ручных проверок и
позволяет раньше реагировать на перегрев, переувлажнение или сбои, это
может уменьшать трудозатраты и риск брака. В тексте ВКР эти эффекты
должны быть представлены как сценарные, с явно указанными исходными
допущениями.

\subsection{3.11 Подготовка
ML-эксперимента}\label{ux43fux43eux434ux433ux43eux442ux43eux432ux43aux430-ml-ux44dux43aux441ux43fux435ux440ux438ux43cux435ux43dux442ux430}

Фактическая история позволяет перейти от декларации ``на основе
машинного обучения'' к воспроизводимому baseline-эксперименту.
Минимальная постановка состоит в прогнозе температуры и влажности на
горизонте 10-30 минут по лаговым значениям датчиков, времени суток и
состояниям реле. CO\(_2\) следует прогнозировать отдельно после
фильтрации технических нулей и проверки качества ряда.

Формирование ML-таблицы включает:

\begin{enumerate}
\def\labelenumi{\arabic{enumi})}
\tightlist
\item
  ресемплирование сенсорных рядов к равномерному шагу;
\item
  разворот устройств в отдельные признаки;
\item
  восстановление состояний реле на временной шкале;
\item
  исключение технических нулей;
\item
  построение лагов 5, 10, 30 и 60 минут;
\item
  добавление признаков часа суток и дня недели;
\item
  сдвиг целевой переменной на горизонт прогноза;
\item
  временное разделение train/test без случайного перемешивания.
\end{enumerate}

В качестве baseline используется прогноз последним наблюдением. С ним
следует сравнить линейную регрессию и одну табличную ML-модель, например
random forest или gradient boosting. Метрики: MAE, RMSE и R2. Успешным
результатом для ВКР будет не внедрение автономного ML-управления, а
доказательство, что собранная платформой история может быть
преобразована в обучающий датасет и что модель хотя бы для части
параметров превосходит naive baseline.

\subsection{3.12 Ограничения опытного
внедрения}\label{ux43eux433ux440ux430ux43dux438ux447ux435ux43dux438ux44f-ux43eux43fux44bux442ux43dux43eux433ux43e-ux432ux43dux435ux434ux440ux435ux43dux438ux44f}

Опытное внедрение подтвердило работоспособность платформы и наличие
реальной истории, но имеет ограничения:

\begin{enumerate}
\def\labelenumi{\arabic{enumi})}
\tightlist
\item
  данные являются наблюдательными, а не результатом специально
  поставленного факторного эксперимента;
\item
  часть событий реле совпадает во времени, поэтому эффекты отдельных
  каналов трудно разделить;
\item
  отсутствует непрерывная наружная температура, наружная влажность и
  солнечная радиация;
\item
  у ряда \texttt{co2-1} есть технические нули и разрывы;
\item
  ручные измерения по салату не всегда имеют точное время, достаточное
  для синхронизации с минутным рядом микроклимата;
\item
  мощности отдельных исполнительных устройств требуют уточнения для
  финального экономического расчета.
\end{enumerate}

Эти ограничения не отменяют результата ВКР. Напротив, они показывают,
какие требования должна закрывать следующая итерация платформы: добавить
наружный климат, фиксировать зоны и партии растений, уточнить паспортные
мощности оборудования, синхронизировать ручные измерения по времени и
заранее планировать controlled experiment.

\subsection{3.13 Выводы по
главе}\label{ux432ux44bux432ux43eux434ux44b-ux43fux43e-ux433ux43bux430ux432ux435-1}

В третьей главе описана фактическая реализация и опытное внедрение
IoT/ML-платформы управления микроклиматом теплицы. Реализованы сенсорные
узлы, релейный контроллер, MQTT-обмен, dashboard, правила, профили,
SQLite-журнал и экспорт данных. С Raspberry Pi получена реальная история
за период с 28 марта по 26 мая 2026 года: 243946 строк сенсорных данных
и 165 событий реле.

Анализ показал, что платформа не только отображает текущие значения, но
и формирует исследовательский датасет. В нем можно анализировать
суточную динамику, расхождение датчиков, качество данных, активность
реле, естественные эксперименты включения оборудования и связь с
агробиологическими данными по салату. Основной результат практической
части состоит в создании измеряемого и управляемого объекта, пригодного
для дальнейшего ML-прогнозирования и расчетной оценки потенциального
экономического эффекта.

\VkrStructuralSection{ЗАКЛЮЧЕНИЕ}\label{ux437ux430ux43aux43bux44eux447ux435ux43dux438ux435}

В ходе работы разработан и опытно внедрен прототип расширяемой
IoT/ML-платформы мониторинга и управления микроклиматом теплицы. Теплица
рассматривалась как теплотехнический и киберфизический объект, в котором
температура, влажность, концентрация CO\(_2\), вентиляция, нагрев,
увлажнение и освещение должны анализироваться совместно с управляющими
воздействиями.

В аналитической части рассмотрены современные подходы к smart
greenhouse, IoT-архитектурам, прогнозированию микроклимата,
MPC/RL-подходам, энергоэффективному управлению и обработке неоднородных
сенсорных данных. Обзор показал, что машинное обучение в тепличных
системах требует надежной инженерной основы: без устойчивого сбора
данных, управления оборудованием и хранения истории невозможно перейти к
корректному прогнозированию и интеллектуальной поддержке решений.

В проектной части сформирована архитектура расширяемой платформы. Она
включает сенсорный слой, исполнительный контроллер, MQTT-брокер,
web-dashboard, правила автоматизации, профили управления, SQLite-журнал
и экспорт данных. Отдельно описаны реестр устройств и метрик,
унифицированная модель MQTT-обмена, перспективная модель данных
эксперимента выращивания, математический каркас тепловлажностного режима
и расчетная модель потенциального экономического эффекта.

В практической части реализованы сенсорные узлы температуры, влажности и
CO\(_2\), релейный контроллер на 15 каналов, MQTT-обмен, web-панель
оператора, правила, профили, журналирование истории и экспорт данных. С
опытного объекта получена фактическая история эксплуатации за период с
28 марта по 26 мая 2026 года: 243946 строки сенсорных данных и 165
событий реле. Это подтверждает, что система не только отображает текущие
параметры, но и формирует воспроизводимый датасет для анализа.

На основе накопленной истории выполнен разведочный анализ микроклимата и
управляющих воздействий. Выявлены различия между датчиками, технические
нули и разрывы данных, активность реле притока и увлажнения, а также
возможность выделения естественных экспериментов включения оборудования.
Такие эксперименты не являются строгим факторным доказательством
причинности, но показывают, что платформа фиксирует микроклимат и
управляющие воздействия на общей временной шкале.

Данные по выращиванию салата связаны с платформенной рамкой через
фотопериод 18/6, PPFD, DLI, электрическую мощность световых вариантов,
EC, биометрию и потери массы при хранении. В работе зафиксировано
ограничение: имеющиеся агробиологические данные можно использовать как
основу экспериментальной рамки, но не как доказательство оптимального
режима выращивания без следующего controlled experiment.

Основные результаты работы:

\begin{enumerate}
\def\labelenumi{\arabic{enumi})}
\tightlist
\item
  Обоснована исследовательская рамка ВКР как разработки расширяемой
  IoT/ML-платформы управления микроклиматом теплицы.
\item
  Сформирован аналитический обзор smart greenhouse, IoT,
  ML-прогнозирования, методов управления и энергоэффективности.
\item
  Спроектирована архитектура платформы с единой моделью обмена данными,
  журналированием датчиков и исполнительных воздействий.
\item
  Реализован программно-аппаратный прототип с сенсорными узлами,
  релейным контроллером, MQTT, dashboard, правилами, профилями и
  SQLite-историей.
\item
  Получен фактический датасет опытной эксплуатации теплицы.
\item
  Подготовлены таблицы и графики разведочного анализа данных.
\item
  Сформулирована постановка ML-эксперимента с baseline-прогнозом,
  лаговыми признаками и метриками MAE, RMSE и R2.
\item
  Описана расчетная методика оценки потенциального экономического
  эффекта через энергозатраты, DLI, работу исполнительных устройств,
  потери продукции и трудозатраты оператора.
\end{enumerate}

Цель работы достигнута: создан прототип платформы, обеспечивающий
подключение разнородных датчиков, управление исполнительными
устройствами, журналирование истории и подготовку данных для машинного
обучения. При этом работа не заявляет разработку нового ML-алгоритма и
не утверждает завершенное автономное ML-управление. Машинное обучение
рассматривается как верхний прогнозный слой, который должен опираться на
безопасный нижний контур правил и профилей.

Основные ограничения связаны с наблюдательным характером данных,
отсутствием непрерывной наружной погоды, техническими нулями части
сенсорных рядов, малым числом событий отдельных реле и недостаточной
синхронизацией ручных измерений по растениям. Эти ограничения задают
направления дальнейшего развития: добавить наружные датчики, уточнить
паспортные мощности оборудования, синхронизировать агробиологические
измерения по времени, выполнить baseline ML-эксперимент и затем провести
controlled experiment с заранее заданными режимами.

Таким образом, выполненная работа создает основу для перехода от ручного
и правилового управления к data-driven подходу, при котором решения по
микроклимату теплицы опираются на историю наблюдений, анализ
естественных экспериментов, прогноз реакции среды и расчетную оценку
энергоэффективности.

\VkrStructuralSection{СПИСОК ИСПОЛЬЗОВАННЫХ ИСТОЧНИКОВ}\label{ux441ux43fux438ux441ux43eux43a-ux438ux441ux43fux43eux43bux44cux437ux43eux432ux430ux43dux43dux44bux445-ux438ux441ux442ux43eux447ux43dux438ux43aux43eux432}

\begin{enumerate}
\def\labelenumi{\arabic{enumi}.}
\tightlist
\item
  El Ouaham W., Sadik M., Ennajih A., Mouzouna Y., Orchi H., Elouaham S.
  Smart Greenhouses in the Era of IoT and AI: A Comprehensive Review of
  AI Applications, Spectral Sensing, Multimodal Data Fusion, and
  Intelligent Systems // Agriculture. 2026. Vol. 16, No.~7. Article 761.
  DOI: 10.3390/agriculture16070761.
\item
  Bersani C., Ruggiero C., Sacile R., Soussi A., Zero E. Internet of
  Things Approaches for Monitoring and Control of Smart Greenhouses in
  Industry 4.0 // Energies. 2022. Vol. 15, No.~10. Article 3834. DOI:
  10.3390/en15103834.
\item
  Ван Вэйянь, Харисов А. Р. Сравнительный анализ современных систем
  автоматизированного управления микроклиматом в теплице // Вестник
  науки. 2025. № 12 (93). С. 256-275. URL:
  \url{https://cyberleninka.ru/article/n/sravnitelnyy-analiz-sovremennyh-sistem-avtomatizirovannogo-upravleniya-mikroklimatom-v-teplitse.}
\item
  Халина Т. М., Цыбанюк Д. Е. Разработка автоматизированной системы
  управления микроклиматом тепличного комплекса // Наука и молодежь:
  материалы XVII Всероссийской научно-технической конференции. Барнаул:
  АлтГТУ, 2020. С. 125-126.
\item
  Тарков Ю. М., Сукьясов С. В. Построение модели управления параметрами
  микроклимата в теплице // Материалы Иркутского ГАУ. 2025. С. 836-839.
\item
  Воробьева Н. С. {[}и др.{]}. Интеллектуальная система контроля
  температуры в теплице // Известия Нижневолжского агроуниверситетского
  комплекса. 2026. № 2 (86). С. 517-522.
\item
  Кельсин А. А. Автоматизированная система управления тепличным
  хозяйством на базе микроконтроллеров: выпускная квалификационная
  работа. Екатеринбург: Уральский государственный педагогический
  университет, 2021.
\item
  Гречиха А. С. Обзор системы управления микроклиматом
  автоматизированной теплицы для выращивания микрозелени // Агротехника
  и энергообеспечение. 2023. № 2.
\item
  Alhnaity B., Pearson S., Leontidis G., Kollias S. Using Deep Learning
  to Predict Plant Growth and Yield in Greenhouse Environments.
  arXiv:1907.00624. 2019. URL: \url{https://arxiv.org/abs/1907.00624.}
\item
  Gong L., Yu M., Jiang S., Cutsuridis V., Pearson S. Deep Learning
  Based Prediction on Greenhouse Crop Yield Combined TCN and RNN //
  Sensors. 2021. Vol. 21, No.~13. Article 4537. DOI: 10.3390/s21134537.
\item
  Ahn J. Y., Kim Y., Park H., Park S. H., Suh H. K. Evaluating
  Time-Series Prediction of Temperature, Relative Humidity, and CO\(_2\)
  in the Greenhouse with Transformer-Based and RNN-Based Models //
  Agronomy. 2024. Vol. 14, No.~3. Article 417. DOI:
  10.3390/agronomy14030417.
\item
  Huang S., Liu Q., Wu Y., Chen M., Yin H., Zhao J. Edible Mushroom
  Greenhouse Environment Prediction Model Based on Attention CNN-LSTM //
  Agronomy. 2024. Vol. 14, No.~3. Article 473. DOI:
  10.3390/agronomy14030473.
\item
  Choi W.-J., Yang M. Probabilistic Deep Learning Framework for
  Greenhouse Microclimate Prediction with Time-Varying Uncertainty and
  Covariance Analysis // Agriculture. 2025. Vol. 15, No.~23. Article
  2461. DOI: 10.3390/agriculture15232461.
\item
  Aborujilah A., Al-Sarem M., Abu-Zanona M. A. Forecast-Driven Climate
  Control for Smart Greenhouses: Energy Optimization Using LSTM Model //
  Energies. 2025. Vol. 18, No.~21. Article 5821. DOI:
  10.3390/en18215821.
\item
  Soumo E. A., Calisti M., Polydoros A. Data-Driven Greenhouse Climate
  Regulation in Lettuce Cultivation Using BiLSTM and GRU Predictive
  Control. arXiv:2507.21669. 2026. URL:
  \url{https://arxiv.org/abs/2507.21669.}
\item
  Thwin K. M. M., Horanont T., Phatrapornnant T. Machine-Learning
  Microclimate Forecasting for Adaptive Equipment Control via Web
  Integration in Open-Ventilated Greenhouses // AgriEngineering. 2024.
  Vol. 6, No.~3. P. 2845-2869. DOI: 10.3390/agriengineering6030165.
\item
  Mallick S., Airaldi F., Dabiri A., Sun C., De Schutter B.
  Reinforcement learning-based model predictive control for greenhouse
  climate control // Smart Agricultural Technology. 2025. Vol. 10.
  Article 100751. DOI: 10.1016/j.atech.2024.100751.
\item
  Msaad S., Harraway M., McAllister R. D. RL-Guided MPC for Autonomous
  Greenhouse Control. arXiv:2506.13278. 2025. URL:
  \url{https://arxiv.org/abs/2506.13278.}
\item
  Morcego B., Yin W., Boersma S., van Henten E., Puig V., Sun C.
  Reinforcement Learning Versus Model Predictive Control on Greenhouse
  Climate Control. arXiv:2303.06110. 2023. URL:
  \url{https://arxiv.org/abs/2303.06110.}
\item
  Xiao M., Lan J., Yu J., Ma W., Xie Q., Sun C. Grower-in-the-Loop
  Interactive Reinforcement Learning for Greenhouse Climate Control.
  arXiv:2505.23355. 2025. URL: \url{https://arxiv.org/abs/2505.23355.}
\item
  Platero-Horcajadas M., Pardo-Pina S., Cámara-Zapata J.-M.,
  Brenes-Carranza J.-A., Ferrández-Pastor F.-J. Enhancing Greenhouse
  Efficiency: Integrating IoT and Reinforcement Learning for Optimized
  Climate Control // Sensors. 2024. Vol. 24, No.~24. Article 8109. DOI:
  10.3390/s24248109.
\item
  Jawad M., Wahid F., Ali S., Ma Y., Alkhyyat A., Khan J., Lee Y. Energy
  optimization and plant comfort management in smart greenhouses using
  the artificial bee colony algorithm // Scientific Reports. 2025. DOI:
  10.1038/s41598-024-84141-5.
\item
  Teja G. S., Sathish P. MQTT Protocol based Smart Greenhouse
  Environment Monitoring System using Machine Learning // International
  Journal of Innovative Technology and Exploring Engineering. 2020. Vol.
  9, No.~9. P. 278-283. DOI: 10.35940/ijitee.I7149.079920.
\item
  Joni I. M. {[}et al.{]}. Smart Farming Experiment: IoT-Enhanced
  Greenhouse Design for Rice Cultivation with Foliar and Soil
  Fertilization // AgriEngineering. 2025. Vol. 7, No.~11. Article 380.
  DOI: 10.3390/agriengineering7110380.
\item
  Wu Y., Fu X. Clean and Smart Energy Technologies for Agricultural
  Energy Internet Systems: A Comprehensive Review and Future
  Perspectives // Applied Sciences. 2026. Vol. 16, No.~10. Article 4859.
  DOI: 10.3390/app16104859.
\item
  Hindi I., Alsharkawi A., Al-Ajlouni M., Qarallah B. Enhancing
  autonomous agriculture control systems in greenhouses for sustainable
  resource usage using deep learning techniques // PLOS One. 2026. DOI:
  10.1371/journal.pone.0344946.
\item
  Тюхтий Ю. А. Разработка интеллектуальной системы управления
  температурой в помещении: магистерская диссертация. Екатеринбург:
  Уральский федеральный университет, 2021.
\item
  Chen S., Liu A., Tang F., Hou P., Lu Y., Yuan P. A Review of
  Environmental Control Strategies and Models for Modern Agricultural
  Greenhouses // Sensors. 2025. Vol. 25, No.~5. Article 1388. DOI:
  10.3390/s25051388.
\item
  Ogunlowo Q. O., Akpenpuun T. D., Na W.-H., Rabiu A., Adesanya M. A.,
  Addae K. S., Kim H.-T., Lee H.-W. Analysis of Heat and Mass
  Distribution in a Single- and Multi-Span Greenhouse Microclimate //
  Agriculture. 2021. Vol. 11, No.~9. Article 891. DOI:
  10.3390/agriculture11090891.
\item
  Li H., Li A., Hou Y., Zhang C., Guo J., Li J., Ma Y., Wang T., Yin Y.
  Analysis of Heat and Humidity in Single-Slope Greenhouses with Natural
  Ventilation // Buildings. 2023. Vol. 13, No.~3. Article 606. DOI:
  10.3390/buildings13030606.
\item
  Ghaderi M., Reddick C., Sorin M. A Systematic Heat Recovery Approach
  for Designing Integrated Heating, Cooling, and Ventilation Systems for
  Greenhouses // Energies. 2023. Vol. 16, No.~14. Article 5493. DOI:
  10.3390/en16145493.
\item
  Самарин Г. Н., Попов А. Н., Плотников С. В., Ружьев В. А. Методика
  построения имитационной модели оптимального управления параметрами
  микроклимата сельскохозяйственного объекта // Аграрный научный журнал.
  2025. № 1. С. 126-135. DOI: 10.28983/asj.y2025i1pp126-135.
\end{enumerate}


\end{document}
